% Pertemuan 9 Praktikum Pemrograman (IF25-11002). Dihasilkan oleh tools/build.py dari tools/konten/p09.py.
% Kompilasi: xelatex P09.tex (dua kali). Catatan: xelatex -jobname=P09-catatan "\def\catatan{1}% Pertemuan 9 Praktikum Pemrograman (IF25-11002). Dihasilkan oleh tools/build.py dari tools/konten/p09.py.
% Kompilasi: xelatex P09.tex (dua kali). Catatan: xelatex -jobname=P09-catatan "\def\catatan{1}% Pertemuan 9 Praktikum Pemrograman (IF25-11002). Dihasilkan oleh tools/build.py dari tools/konten/p09.py.
% Kompilasi: xelatex P09.tex (dua kali). Catatan: xelatex -jobname=P09-catatan "\def\catatan{1}% Pertemuan 9 Praktikum Pemrograman (IF25-11002). Dihasilkan oleh tools/build.py dari tools/konten/p09.py.
% Kompilasi: xelatex P09.tex (dua kali). Catatan: xelatex -jobname=P09-catatan "\def\catatan{1}\input{P09.tex}"
\documentclass[t]{beamer}
\usetheme{ITERA}
\input{catatan-preamble.tex}
\renewcommand\ITERAmeeting{Pertemuan 9}
\begin{document}
\SlideJudul{IF25-11002 PRAKTIKUM PEMROGRAMAN}{Pertemuan 9\\Array Satu Dimensi}{Menyimpan banyak nilai sejenis, mengaksesnya dengan indeks, dan memprosesnya dengan perulangan}{Tim Pengajar}{Tahun 2026. Sejajar dengan Pertemuan 9 Algoritma Pemrograman: Array 1D.}
\note{Buka dengan menanyakan satu hal dari pertemuan lalu yang masih membingungkan.}

\begin{frame}{Dua mata kuliah, satu minggu}
\framesubtitle{Sebelum mulai}
Topik minggu ini sama di dua mata kuliah, dengan peran yang berbeda.
\vspace{10pt}

\begin{columns}[T]
\begin{column}{0.47\textwidth}
\begin{Blok}{Algoritma: Array 1D}
Konsep larik, indeks, dan algoritma dasar pada larik: penjumlahan, maksimum, dan pembalikan.
\end{Blok}
\end{column}
\begin{column}{0.47\textwidth}
\begin{BlokGelap}{Praktikum: Array 1D}
Mendeklarasikan array C++, mengisi dan membacanya dengan \texttt{for}, serta menghindari akses di luar batas.
\end{BlokGelap}
\end{column}
\end{columns}
\note{Tanyakan apa yang sudah dibahas di Algoritma minggu ini. Gunakan jawabannya sebagai jembatan ke praktikum.}
\end{frame}

\begin{frame}{Saat pulang nanti, kalian bisa}
\framesubtitle{Target hari ini}
\begin{columns}[T]
\begin{column}{0.47\textwidth}
\begin{Langkah}{1}{Menerapkan}
Menulis program yang menyimpan data sejenis di array satu dimensi dan memprosesnya dengan perulangan: mengisi, menampilkan, menjumlahkan, mencari maksimum, dan menggeser

{\footnotesize\color{ITERAInkMuted}Sub-CPMK 9.1, CPMK0607}
\end{Langkah}
\end{column}
\begin{column}{0.47\textwidth}
\begin{Langkah}{2}{Menjelaskan}
Menelusuri isi array sesudah setiap langkah pemrosesan dan menjelaskan kesalahan indeks, termasuk akses di luar batas

{\footnotesize\color{ITERAInkMuted}Sub-CPMK 9.2, CPMK0608}
\end{Langkah}
\end{column}
\end{columns}
\vspace{8pt}
\begin{Catatan}
Dua kemampuan ini dinilai sama berat di Exit Ticket, tugas, UTS, dan UAS.
\end{Catatan}
\note{Bacakan target dengan pelan. Kata kuncinya menerapkan dan menjelaskan.}
\end{frame}

\begin{frame}{Ingat minggu lalu}
\framesubtitle{Review, 5 menit}
\begin{columns}[T]
\begin{column}{0.31\textwidth}
\begin{BlokGelap}{Pertanyaan 1}
Apa hasil \texttt{int n = 4096} setelah dijumlah digitnya?
\end{BlokGelap}
\end{column}
\begin{column}{0.31\textwidth}
\begin{BlokGelap}{Pertanyaan 2}
Berapa kali badan \texttt{for i 1..3, for j 1..i} dijalankan?
\end{BlokGelap}
\end{column}
\begin{column}{0.31\textwidth}
\begin{BlokGelap}{Pertanyaan 3}
Apa beda soal A dan soal B di UTS?
\end{BlokGelap}
\end{column}
\end{columns}
\vspace{8pt}
Jawab di kertas dulu, lalu cocokkan dengan teman sebelah.\note{Pertemuan pertama sesudah UTS. Bahas singkat kesalahan umum UTS sebelum masuk materi.}
\end{frame}

\Bagian{1}{Masalah pembuka}{Siapa yang di atas rata-rata?}
\note{Masalah pembuka 10 menit.}

\begin{frame}[fragile]{Siapa yang di atas rata-rata?}
\framesubtitle{Masalah minggu ini}
\begin{columns}[T]
\begin{column}{0.55\textwidth}
{\small Dosen ingin tahu berapa mahasiswa yang nilainya di atas rata-rata kelas.

\vspace{6pt}
Dengan pola sentinel minggu 5, rata-rata baru diketahui setelah semua nilai dibaca. Tetapi saat itu nilai-nilai sebelumnya sudah hilang, karena setiap \texttt{cin} menimpa variabel yang sama.\par}
\vspace{6pt}
\begin{Catatan}
\textbf{Diskusi 2 menit.} Apa yang perlu dilakukan supaya semua nilai masih ada saat rata-rata sudah diketahui? Berapa variabel yang dibutuhkan untuk 40 mahasiswa?
\end{Catatan}
\end{column}
\begin{column}{0.41\textwidth}
\begin{Keluaran}[]
Nilai: 70 85 60 90 75 80
Rata-rata   : 76.67
Di atas rata: 3
\end{Keluaran}
\end{column}
\end{columns}
\note{Tampung beberapa jawaban tanpa menilai benar salah.}
\end{frame}

\begin{frame}{Dari langkah ke instruksi}
\framesubtitle{Sambungan dengan Algoritma}
\begin{columns}[T]
\begin{column}{0.31\textwidth}
\begin{Langkah}{1}{Simpan semua}
Semua nilai perlu disimpan, bukan hanya totalnya.
\end{Langkah}
\end{column}
\begin{column}{0.31\textwidth}
\begin{Langkah}{2}{Proses dua kali}
Putaran pertama menghitung rata-rata, putaran kedua membandingkan.
\end{Langkah}
\end{column}
\begin{column}{0.31\textwidth}
\begin{Langkah}{3}{Terjemahkan}
Satu array, dua \texttt{for} dengan indeks 0 sampai n − 1.
\end{Langkah}
\end{column}
\end{columns}
\vspace{10pt}
Langkah 1 dan 2 dilatih di Algoritma. Langkah 3 kita kerjakan hari ini dengan C++.\note{Hubungkan eksplisit ke Algoritma minggu ini.}
\end{frame}

\Bagian{2}{Coba dulu, baru dijelaskan}{25 menit eksperimen di GDB Online}
\note{Struggle 25 menit. Berkeliling, jangan menjelaskan materi dulu.}

\begin{frame}{Misi kalian}
\framesubtitle{Struggle, 25 menit}
\begin{columns}[T]
\begin{column}{0.31\textwidth}
\begin{Langkah}{1}{Deklarasi}
Buat \texttt{int a[5] = \{70, 85, 60, 90, 75\};} dan tampilkan \texttt{a[0]} dan \texttt{a[4]}.
\end{Langkah}
\end{column}
\begin{column}{0.31\textwidth}
\begin{Langkah}{2}{Ulangi}
Tampilkan semua elemen dengan \texttt{for}.
\end{Langkah}
\end{column}
\begin{column}{0.31\textwidth}
\begin{Langkah}{3}{Uji batas}
Tampilkan \texttt{a[5]}. Apa yang muncul?
\end{Langkah}
\end{column}
\end{columns}
\vspace{8pt}
\begin{columns}[T]
\begin{column}{0.47\textwidth}
\begin{Blok}{Boleh}
Bertanya ke teman, mengubah apa saja, sengaja membuat error.
\end{Blok}
\end{column}
\begin{column}{0.47\textwidth}
\begin{BlokAlert}{Belum dulu}
Mencari jawaban jadi di internet atau AI.
\end{BlokAlert}
\end{column}
\end{columns}
\note{Catat kesalahan yang sering muncul untuk dibahas di debrief.}
\end{frame}

\begin{frame}{Apa yang kalian temukan?}
\framesubtitle{Debrief struggle}
\begin{columns}[T]
\begin{column}{0.31\textwidth}
\begin{BlokGelap}{Pertanyaan 1}
Mengapa indeks pertama 0, bukan 1?
\end{BlokGelap}
\end{column}
\begin{column}{0.31\textwidth}
\begin{BlokGelap}{Pertanyaan 2}
Apa yang tampil untuk \texttt{a[5]}, dan apakah sama di setiap komputer?
\end{BlokGelap}
\end{column}
\begin{column}{0.31\textwidth}
\begin{BlokGelap}{Pertanyaan 3}
Bagaimana kalian menghitung jumlah semua elemen?
\end{BlokGelap}
\end{column}
\end{columns}
\vspace{10pt}
Jawaban kalian akan kita uji di bagian berikutnya.\note{Tulis jawaban di papan; buktikan atau patahkan di bagian materi.}
\end{frame}

\Bagian{3}{Banyak nilai, satu nama}{Deklarasi, indeks, perulangan, dan algoritma dasar array}
\note{Materi 40 menit. Tunjukkan setiap contoh langsung di GDB Online.}

\begin{frame}[fragile]{Deklarasi dan indeks}
\framesubtitle{Array satu dimensi}
\begin{columns}[T]
\begin{column}{0.55\textwidth}
\begin{Kode}[]{array.cpp}
int nilai[5] = {70, 85, 60, 90, 75};
cout << nilai[0] << " " << nilai[4] << endl;
nilai[2] = 65;
nilai[3] += 5;
for (int i = 0; i < 5; i++) {
    cout << nilai[i] << " ";
}
cout << endl;
\end{Kode}
\end{column}
\begin{column}{0.41\textwidth}
\only<1>{\begin{TebakDulu}
{\ITERAKickerFont\fontsize{10pt}{12pt}\selectfont\color{ITERAAlert}TEBAK DULU}\par\vspace{4pt}
Sebelum dijalankan, tulis keluarannya di kertas.
\end{TebakDulu}
}\only<2>{\KeluaranFile[]{keluaran/P09/k001.txt}
\begin{Catatan}
Array berukuran 5 punya indeks 0 sampai 4. Ukuran harus konstanta.
\end{Catatan}
}
\end{column}
\end{columns}
\note<1>{Minta mahasiswa menebak keluaran sebelum membuka slide berikutnya. Tanyakan satu atau dua tebakan yang berbeda, lalu buktikan.}
\end{frame}

\begin{frame}[fragile]{Mengisi array dari masukan}
\framesubtitle{Pola baca-proses}
\begin{columns}[T]
\begin{column}{0.55\textwidth}
\begin{Kode}[]{baca.cpp}
int n, a[100];
cin >> n;
for (int i = 0; i < n; i++) {
    cin >> a[i];
}
int total = 0;
for (int i = 0; i < n; i++) {
    total += a[i];
}
cout << "Total " << total << endl;
\end{Kode}
\end{column}
\begin{column}{0.41\textwidth}
\begin{Keluaran}[title={Masukan}]
4
10 20 30 40
\end{Keluaran}
\only<1>{\begin{TebakDulu}
{\ITERAKickerFont\fontsize{10pt}{12pt}\selectfont\color{ITERAAlert}TEBAK DULU}\par\vspace{4pt}
Sebelum dijalankan, tulis keluarannya di kertas.
\end{TebakDulu}
}\only<2>{\KeluaranFile[]{keluaran/P09/k002.txt}
\begin{Catatan}
Deklarasikan array cukup besar (misalnya 100), lalu pakai hanya n elemen pertama.
\end{Catatan}
}
\end{column}
\end{columns}
\note<1>{Minta mahasiswa menebak keluaran sebelum membuka slide berikutnya. Tanyakan satu atau dua tebakan yang berbeda, lalu buktikan.}
\end{frame}

\begin{frame}[fragile]{Mencari maksimum dan posisinya}
\framesubtitle{Algoritma dasar}
\begin{columns}[T]
\begin{column}{0.60\textwidth}
\begin{Kode}[listing options={style=ITERACodeKecil}]{maks.cpp}
int a[6] = {70, 85, 60, 90, 75, 80};
int iMaks = 0;
for (int i = 1; i < 6; i++) {
    if (a[i] > a[iMaks]) {
        iMaks = i;
    }
}
cout << a[iMaks] << " di indeks " << iMaks << endl;
\end{Kode}
\end{column}
\begin{column}{0.36\textwidth}
\only<1>{\begin{TebakDulu}
{\ITERAKickerFont\fontsize{10pt}{12pt}\selectfont\color{ITERAAlert}TEBAK DULU}\par\vspace{4pt}
Sebelum dijalankan, tulis keluarannya di kertas.
\end{TebakDulu}
}\only<2>{\KeluaranFile[listing options={style=ITERAPlainKecil}]{keluaran/P09/k003.txt}
\begin{Catatan}
Simpan indeks, bukan nilainya. Dari indeks, nilai selalu bisa diambil kembali.
\end{Catatan}
}
\end{column}
\end{columns}
\note<1>{Minta mahasiswa menebak keluaran sebelum membuka slide berikutnya. Tanyakan satu atau dua tebakan yang berbeda, lalu buktikan.}
\end{frame}

\begin{frame}[fragile]{Dua putaran: rata-rata lalu bandingkan}
\framesubtitle{Mengapa perlu array}
\begin{columns}[T]
\begin{column}{0.55\textwidth}
\begin{Kode}[]{diatas.cpp}
int a[6] = {70, 85, 60, 90, 75, 80};
int total = 0;
for (int i = 0; i < 6; i++) total += a[i];
double rata = total / 6.0;
int diAtas = 0;
for (int i = 0; i < 6; i++) {
    if (a[i] > rata) diAtas++;
}
cout << rata << " " << diAtas << endl;
\end{Kode}
\end{column}
\begin{column}{0.41\textwidth}
\only<1>{\begin{TebakDulu}
{\ITERAKickerFont\fontsize{10pt}{12pt}\selectfont\color{ITERAAlert}TEBAK DULU}\par\vspace{4pt}
Sebelum dijalankan, tulis keluarannya di kertas.
\end{TebakDulu}
}\only<2>{\KeluaranFile[]{keluaran/P09/k004.txt}
}
\end{column}
\end{columns}
\note<1>{Minta mahasiswa menebak keluaran sebelum membuka slide berikutnya. Tanyakan satu atau dua tebakan yang berbeda, lalu buktikan.}
\end{frame}

\begin{frame}[fragile]{Di luar batas}
\framesubtitle{Kesalahan paling berbahaya}
\begin{columns}[T]
\begin{column}{0.55\textwidth}
\begin{Kode}[]{batas.cpp}
int a[3] = {1, 2, 3};
int jumlah = 0;
for (int i = 0; i < 3; i++) {
    jumlah += a[i];
}
cout << jumlah << endl;
// for (int i = 0; i <= 3; i++) membaca a[3]:
// di luar array, hasilnya tidak bisa ditebak
\end{Kode}
\end{column}
\begin{column}{0.41\textwidth}
\only<1>{\begin{TebakDulu}
{\ITERAKickerFont\fontsize{10pt}{12pt}\selectfont\color{ITERAAlert}TEBAK DULU}\par\vspace{4pt}
Sebelum dijalankan, tulis keluarannya di kertas.
\end{TebakDulu}
}\only<2>{\KeluaranFile[]{keluaran/P09/k005.txt}
\begin{Catatan}
C++ tidak memeriksa batas indeks. Program bisa tetap jalan dengan hasil acak, atau berhenti mendadak.
\end{Catatan}
}
\end{column}
\end{columns}
\note<1>{Minta mahasiswa menebak keluaran sebelum membuka slide berikutnya. Tanyakan satu atau dua tebakan yang berbeda, lalu buktikan.}
\end{frame}

\begin{frame}{Menggeser dan membalik}
\framesubtitle{Menelusuri isi array}
Rotasi kanan \{10, 20, 30, 40, 50\}: simpan elemen terakhir, geser dari belakang, isi depan.
\vspace{8pt}

\small
\begin{tabular}{@{}>{\raggedright\arraybackslash}p{172pt}>{\raggedright\arraybackslash}p{364pt}>{\raggedright\arraybackslash}p{326pt}@{}}
\kepala{Langkah} & \kepala{Isi array} & \kepala{Keterangan}\\
awal & 10 20 30 40 50 & simpan = 50\\
\barisgenap i = 4 & 10 20 30 40 40 & a[4] = a[3]\\
i = 3 & 10 20 30 30 40 & a[3] = a[2]\\
\barisgenap i = 1 & 10 10 20 30 40 & setelah i = 2 dan i = 1\\
akhir & 50 10 20 30 40 & a[0] = simpan\\
\garisdasar
\end{tabular}
\end{frame}

\begin{frame}[fragile]{Tebak keluarannya}
\framesubtitle{Latihan menjelaskan, 3 menit}
\begin{columns}[T]
\begin{column}{0.56\textwidth}
\begin{Kode}[]{tebak.cpp}
int a[5] = {3, 1, 4, 1, 5};
for (int i = 1; i < 5; i++) {
    a[i] = a[i] + a[i - 1];
}
for (int i = 0; i < 5; i++) {
    cout << a[i] << " ";
}
cout << endl;
\end{Kode}
\end{column}
\begin{column}{0.40\textwidth}
\begin{Catatan}
Tulis keluarannya di kertas, karakter demi karakter. Jangan dijalankan dulu.
\end{Catatan}
\end{column}
\end{columns}
\note{Contoh soal Bagian B. Beri waktu 3 menit sebelum slide jawaban.}
\end{frame}

\begin{frame}[fragile]{Tebak keluarannya: jawaban}
\framesubtitle{Latihan menjelaskan}
\begin{columns}[T]
\begin{column}{0.56\textwidth}
\begin{Kode}[]{tebak.cpp}
int a[5] = {3, 1, 4, 1, 5};
for (int i = 1; i < 5; i++) {
    a[i] = a[i] + a[i - 1];
}
for (int i = 0; i < 5; i++) {
    cout << a[i] << " ";
}
cout << endl;
\end{Kode}
\end{column}
\begin{column}{0.40\textwidth}
\begin{Keluaran}[]
3 4 8 9 14 
\end{Keluaran}
\begin{Catatan}
Setiap elemen ditambah elemen sebelumnya yang sudah diperbarui, sehingga terbentuk jumlah kumulatif: 3, 4, 8, 9, 14.
\end{Catatan}
\end{column}
\end{columns}
\note{Tanyakan jawaban yang paling sering salah, lalu telusuri bersama.}
\end{frame}

\begin{frame}{Error yang sering muncul minggu ini}
\framesubtitle{Pesan diambil langsung dari g++}
\footnotesize
\begin{tabular}{@{}>{\raggedright\arraybackslash}p{208pt}>{\raggedright\arraybackslash}p{256pt}>{\raggedright\arraybackslash}p{398pt}@{}}
\kepala{Kesalahan} & \kepala{Contoh} & \kepala{Pesan pertama dari g++}\\
Ukuran negatif & \texttt{int a[-3];} & \texttt{error: narrowing conversion of '-3' from 'int' to 'long unsigned int'}\\
\barisgenap Inisialisasi kelebihan elemen & \texttt{int a[2] = \{1, 2, 3\};} & \texttt{error: too many initializers for 'int [2]'}\\
Indeks bukan bilangan bulat & \texttt{a[1.5]} & \texttt{error: invalid types 'int [3][double]' for array subscript}\\
\barisgenap Menugaskan array ke array & \texttt{b = a;} & \texttt{error: invalid array assignment}\\
Ukuran tidak ditulis & \texttt{int a[];} & \texttt{error: storage size of 'a' isn't known}\\
\garisdasar
\end{tabular}
\vspace{10pt}

{\small Pesan di tabel ini dihasilkan dengan g++ dan opsi -Wall, sama seperti di cek.sh.}
\note{Minta mahasiswa memotret slide ini.}
\end{frame}

\Bagian{4}{Hands-on bertingkat}{45 menit, dari dasar sampai tantangan}
\note{Mahasiswa membuka folder 02 Hands-on/P9 - Array 1D - Hands-on.}

\begin{frame}{Peta hands-on}
\framesubtitle{Folder 02 Hands-on/P9 - Array 1D - Hands-on}
\small
\begin{tabular}{@{}>{\raggedright\arraybackslash}p{36pt}>{\raggedright\arraybackslash}p{267pt}>{\raggedright\arraybackslash}p{110pt}>{\raggedright\arraybackslash}p{83pt}>{\raggedright\arraybackslash}p{341pt}@{}}
\kepala{No} & \kepala{File} & \kepala{Tingkat} & \kepala{Waktu} & \kepala{Yang dilatih}\\
1 & \texttt{handson1-nilai.cpp} & Dasar & 10' & mengisi dan menampilkan, urutan terbalik\\
\barisgenap 2 & \texttt{handson2-statistik.cpp} & Menengah & 15' & rata-rata, maksimum, dua putaran\\
3 & \texttt{handson3-geser.cpp} & Tantangan & 20' & telusuri isi array, perbaiki tiga kesalahan indeks\\
\barisgenap + & \texttt{bonus-frekuensi.cpp} & Pengayaan & sisa & array sebagai pencacah\\
\garisdasar
\end{tabular}
\vspace{10pt}

Kerjakan berurutan. Cocokkan keluaran dengan folder \texttt{expected/}; latihan yang membaca masukan memakai data di folder \texttt{input/}.\note{Saat berkeliling, minta mahasiswa yang mengerjakan latihan tantangan menjelaskan blok PENJELASAN mereka.}
\end{frame}

\begin{frame}{Cara mengecek dan aturannya}
\framesubtitle{Hands-on}
\begin{columns}[T]
\begin{column}{0.47\textwidth}
\begin{Blok}{Di GDB Online}
Salin isi file latihan, lengkapi TODO, klik Run. Bila program membaca masukan, ketik isi file di \texttt{input/}. Bandingkan dengan \texttt{expected/}.
\end{Blok}
\end{column}
\begin{column}{0.47\textwidth}
\begin{BlokGelap}{Di komputer sendiri}
Jalankan \texttt{bash cek.sh}. Skrip mengompilasi dengan \texttt{-Wall -Werror}, memberi masukan dari \texttt{input/}, lalu mencocokkan keluaran.
\end{BlokGelap}
\end{column}
\end{columns}
\vspace{6pt}
\begin{Catatan}
AI boleh untuk memahami pesan error, tidak untuk meminta solusi utuh. Exit Ticket dikerjakan tanpa bantuan apa pun.
\end{Catatan}
\note{Hands-on tidak dinilai langsung; yang dinilai Exit Ticket dan tugas.}
\end{frame}

\Bagian{5}{Exit Ticket dan tugas}{25 menit terakhir, lalu pekerjaan rumah}
\note{Mulai Exit Ticket tepat waktu.}

\begin{frame}{Exit Ticket 9}
\framesubtitle{Individual, 25 menit, tanpa AI dan internet}
\small
\begin{tabular}{@{}>{\raggedright\arraybackslash}p{60pt}>{\raggedright\arraybackslash}p{560pt}>{\raggedright\arraybackslash}p{80pt}>{\raggedright\arraybackslash}p{150pt}@{}}
\kepala{Soal} & \kepala{Isi} & \kepala{Poin} & \kepala{CPMK}\\
A1 & Tulis program yang membaca n nilai ke array dan menampilkan nilai minimum beserta indeksnya & 25 & CPMK0607\\
\barisgenap A2 & Tulis program yang menghitung banyak elemen genap dan ganjil & 25 & CPMK0607\\
B1 & Telusuri isi array setelah sebuah perulangan yang mengubah elemennya & 20 & CPMK0608\\
\barisgenap B2 & Jelaskan mengapa sebuah perulangan menyebabkan akses di luar batas & 20 & CPMK0608\\
B3 & Jelaskan mengapa sebuah masalah butuh array, bukan variabel tunggal & 10 & CPMK0608\\
\garisdasar
\end{tabular}
\vspace{10pt}

{\small Soal A dinilai dari keluaran yang tepat sama. Soal B dinilai dari ketepatan dan alasan. Pelanggaran aturan membuat nilai Exit Ticket ini 0.}
\note{Soal lengkap dibagikan terpisah dan tidak disimpan di repo publik.}
\end{frame}

\begin{frame}{Tugas 9: Sistem Nilai Mahasiswa}
\framesubtitle{Dikumpulkan sebelum Pertemuan 10}
\begin{columns}[T]
\begin{column}{0.47\textwidth}
\begin{Blok}{Bagian A, program, 50 poin}
Buat program dengan ketentuan berikut. Gunakan \texttt{const} untuk ukuran array. Ketentuannya di slide berikut.
\end{Blok}
\end{column}
\begin{column}{0.47\textwidth}
\begin{BlokGelap}{Bagian B, penjelasan, 50 poin}
Komentar maksimal 150 kata dan gambar isi kedua array setelah membaca tiga data contoh. Jelaskan cara menjaga indeks tetap di dalam batas, dan mengapa dua array harus diproses dengan indeks yang sama.
\end{BlokGelap}
\end{column}
\end{columns}
\vspace{8pt}
{\small Data di tugas adalah data kalian sendiri. Rubrik lengkap ada di Modul Belajar 9.}
\note{Ingatkan bahwa penjelasan disalin dari AI mudah dikenali karena tidak menyebut pengalaman mereka sendiri.}
\end{frame}

\begin{frame}{Ketentuan Tugas 9}
\framesubtitle{Sistem Nilai Mahasiswa}
\small
\begin{tabular}{@{}>{\raggedright\arraybackslash}p{87pt}>{\raggedright\arraybackslash}p{788pt}@{}}
\kepala{No} & \kepala{Ketentuan Bagian A}\\
1 & Baca jumlah mahasiswa (maksimal 50)\\
\barisgenap 2 & Baca NIM dan nilai (0 sampai 100) setiap mahasiswa ke dua array sejajar, dengan validasi\\
3 & Tampilkan semua data, nilai tertinggi dan terendah, serta rata-rata kelas\\
\barisgenap 4 & Tampilkan jumlah LULUS (nilai 60 ke atas), TIDAK LULUS, dan persentase kelulusan\\
5 & Fitur pencarian: baca sebuah NIM, tampilkan nilainya atau pesan tidak ditemukan\\
\barisgenap Bonus & Urutkan nilai dari tertinggi ke terendah (materi sorting di Pertemuan 12).\\
\garisdasar
\end{tabular}
\note{Bacakan ketentuan satu per satu dan tanyakan apakah ada yang belum jelas.}
\end{frame}

\begin{frame}{Rubrik Tugas 9}
\framesubtitle{Empat level, sama di semua instrumen}
\footnotesize
\begin{tabular}{@{}>{\raggedright\arraybackslash}p{166pt}>{\raggedright\arraybackslash}p{44pt}>{\raggedright\arraybackslash}p{166pt}>{\raggedright\arraybackslash}p{156pt}>{\raggedright\arraybackslash}p{146pt}>{\raggedright\arraybackslash}p{146pt}@{}}
\kepala{Kriteria} & \kepala{Poin} & \kepala{Sangat baik 85--100} & \kepala{Baik 70--84} & \kepala{Cukup 55--69} & \kepala{Kurang <55}\\
A1 Program berjalan & 20 & tanpa error dan peringatan & ada peringatan & perlu satu perbaikan & tidak terkompilasi\\
\barisgenap A2 Kelengkapan fitur & 15 & semua statistik dan pencarian & satu fitur kurang & dua kurang & tiga atau lebih kurang\\
A3 Validasi dan ketepatan & 15 & validasi tepat, hasil benar untuk n = 1 dan n = 10 & satu salah & dua salah & sebagian besar salah\\
\barisgenap B1 Isi array dan batas indeks & 30 & gambar tepat, alasan jelas & satu langkah keliru & gambar tanpa alasan & tidak ada atau disalin\\
B2 Cerita error dan pengujian & 20 & pesan, sebab, perbaikan, uji & pesan tidak dikutip & hanya perbaikan & tidak ada\\
\garisdasar
\end{tabular}
\vspace{8pt}

{\small Nilai kriteria = poin × skor level. Bagian A total 50, Bagian B total 50.}
\note{Rubrik ini sama strukturnya setiap minggu; yang berubah hanya isi kriteria A2, A3, dan B1.}
\end{frame}

\begin{frame}{Sebelum Pertemuan 10}
\framesubtitle{Persiapan}
\begin{columns}[T]
\begin{column}{0.31\textwidth}
\begin{Langkah}{1}{Selesaikan}
Hands-on yang belum lulus \texttt{cek.sh}, terutama latihan tantangan.
\end{Langkah}
\end{column}
\begin{column}{0.31\textwidth}
\begin{Langkah}{2}{Kumpulkan}
Tugas 9 sebelum Pertemuan 10 dimulai.
\end{Langkah}
\end{column}
\begin{column}{0.31\textwidth}
\begin{Langkah}{3}{Baca}
Modul Belajar 9 untuk mengulang, lalu bagian awal modul berikutnya.
\end{Langkah}
\end{column}
\end{columns}
\vspace{10pt}
Pertemuan 10 membahas array dua dimensi dan string, sejalan dengan Algoritma Pertemuan 10.\note{Tanyakan satu hal yang paling membingungkan hari ini untuk dibuka di review minggu depan.}
\end{frame}

\SlidePenutup{Terima kasih}{Sampai jumpa di Pertemuan 10: array dua dimensi dan string}
\note{Slide penutup.}
\end{document}
"
\documentclass[t]{beamer}
\usetheme{ITERA}
% Mode catatan pembicara: aktif bila \catatan didefinisikan sebelum \input.
\ifdefined\catatan
  \setbeameroption{show only notes}
  \setbeamertemplate{note page}{\vspace*{20pt}\hspace*{4pt}\begin{minipage}{900pt}
    {\ITERAKickerFont\fontsize{11pt}{14pt}\selectfont\color{ITERAGoldText}CATATAN PEMBICARA \ \ |\ \ SLIDE \insertframenumber}\par\vspace{4pt}
    {\ITERADisplay\bfseries\fontsize{22pt}{28pt}\selectfont\color{ITERAStructure}\insertframetitle}\par\vspace{12pt}
    \fontsize{15pt}{23pt}\selectfont\insertnote\end{minipage}}
\fi
\tikzset{fase/.style={draw=none,fill=#1,rounded corners=3pt,minimum width=158pt,minimum height=118pt,text width=146pt,align=center}}

\renewcommand\ITERAmeeting{Pertemuan 9}
\begin{document}
\SlideJudul{IF25-11002 PRAKTIKUM PEMROGRAMAN}{Pertemuan 9\\Array Satu Dimensi}{Menyimpan banyak nilai sejenis, mengaksesnya dengan indeks, dan memprosesnya dengan perulangan}{Tim Pengajar}{Tahun 2026. Sejajar dengan Pertemuan 9 Algoritma Pemrograman: Array 1D.}
\note{Buka dengan menanyakan satu hal dari pertemuan lalu yang masih membingungkan.}

\begin{frame}{Dua mata kuliah, satu minggu}
\framesubtitle{Sebelum mulai}
Topik minggu ini sama di dua mata kuliah, dengan peran yang berbeda.
\vspace{10pt}

\begin{columns}[T]
\begin{column}{0.47\textwidth}
\begin{Blok}{Algoritma: Array 1D}
Konsep larik, indeks, dan algoritma dasar pada larik: penjumlahan, maksimum, dan pembalikan.
\end{Blok}
\end{column}
\begin{column}{0.47\textwidth}
\begin{BlokGelap}{Praktikum: Array 1D}
Mendeklarasikan array C++, mengisi dan membacanya dengan \texttt{for}, serta menghindari akses di luar batas.
\end{BlokGelap}
\end{column}
\end{columns}
\note{Tanyakan apa yang sudah dibahas di Algoritma minggu ini. Gunakan jawabannya sebagai jembatan ke praktikum.}
\end{frame}

\begin{frame}{Saat pulang nanti, kalian bisa}
\framesubtitle{Target hari ini}
\begin{columns}[T]
\begin{column}{0.47\textwidth}
\begin{Langkah}{1}{Menerapkan}
Menulis program yang menyimpan data sejenis di array satu dimensi dan memprosesnya dengan perulangan: mengisi, menampilkan, menjumlahkan, mencari maksimum, dan menggeser

{\footnotesize\color{ITERAInkMuted}Sub-CPMK 9.1, CPMK0607}
\end{Langkah}
\end{column}
\begin{column}{0.47\textwidth}
\begin{Langkah}{2}{Menjelaskan}
Menelusuri isi array sesudah setiap langkah pemrosesan dan menjelaskan kesalahan indeks, termasuk akses di luar batas

{\footnotesize\color{ITERAInkMuted}Sub-CPMK 9.2, CPMK0608}
\end{Langkah}
\end{column}
\end{columns}
\vspace{8pt}
\begin{Catatan}
Dua kemampuan ini dinilai sama berat di Exit Ticket, tugas, UTS, dan UAS.
\end{Catatan}
\note{Bacakan target dengan pelan. Kata kuncinya menerapkan dan menjelaskan.}
\end{frame}

\begin{frame}{Ingat minggu lalu}
\framesubtitle{Review, 5 menit}
\begin{columns}[T]
\begin{column}{0.31\textwidth}
\begin{BlokGelap}{Pertanyaan 1}
Apa hasil \texttt{int n = 4096} setelah dijumlah digitnya?
\end{BlokGelap}
\end{column}
\begin{column}{0.31\textwidth}
\begin{BlokGelap}{Pertanyaan 2}
Berapa kali badan \texttt{for i 1..3, for j 1..i} dijalankan?
\end{BlokGelap}
\end{column}
\begin{column}{0.31\textwidth}
\begin{BlokGelap}{Pertanyaan 3}
Apa beda soal A dan soal B di UTS?
\end{BlokGelap}
\end{column}
\end{columns}
\vspace{8pt}
Jawab di kertas dulu, lalu cocokkan dengan teman sebelah.\note{Pertemuan pertama sesudah UTS. Bahas singkat kesalahan umum UTS sebelum masuk materi.}
\end{frame}

\Bagian{1}{Masalah pembuka}{Siapa yang di atas rata-rata?}
\note{Masalah pembuka 10 menit.}

\begin{frame}[fragile]{Siapa yang di atas rata-rata?}
\framesubtitle{Masalah minggu ini}
\begin{columns}[T]
\begin{column}{0.55\textwidth}
{\small Dosen ingin tahu berapa mahasiswa yang nilainya di atas rata-rata kelas.

\vspace{6pt}
Dengan pola sentinel minggu 5, rata-rata baru diketahui setelah semua nilai dibaca. Tetapi saat itu nilai-nilai sebelumnya sudah hilang, karena setiap \texttt{cin} menimpa variabel yang sama.\par}
\vspace{6pt}
\begin{Catatan}
\textbf{Diskusi 2 menit.} Apa yang perlu dilakukan supaya semua nilai masih ada saat rata-rata sudah diketahui? Berapa variabel yang dibutuhkan untuk 40 mahasiswa?
\end{Catatan}
\end{column}
\begin{column}{0.41\textwidth}
\begin{Keluaran}[]
Nilai: 70 85 60 90 75 80
Rata-rata   : 76.67
Di atas rata: 3
\end{Keluaran}
\end{column}
\end{columns}
\note{Tampung beberapa jawaban tanpa menilai benar salah.}
\end{frame}

\begin{frame}{Dari langkah ke instruksi}
\framesubtitle{Sambungan dengan Algoritma}
\begin{columns}[T]
\begin{column}{0.31\textwidth}
\begin{Langkah}{1}{Simpan semua}
Semua nilai perlu disimpan, bukan hanya totalnya.
\end{Langkah}
\end{column}
\begin{column}{0.31\textwidth}
\begin{Langkah}{2}{Proses dua kali}
Putaran pertama menghitung rata-rata, putaran kedua membandingkan.
\end{Langkah}
\end{column}
\begin{column}{0.31\textwidth}
\begin{Langkah}{3}{Terjemahkan}
Satu array, dua \texttt{for} dengan indeks 0 sampai n − 1.
\end{Langkah}
\end{column}
\end{columns}
\vspace{10pt}
Langkah 1 dan 2 dilatih di Algoritma. Langkah 3 kita kerjakan hari ini dengan C++.\note{Hubungkan eksplisit ke Algoritma minggu ini.}
\end{frame}

\Bagian{2}{Coba dulu, baru dijelaskan}{25 menit eksperimen di GDB Online}
\note{Struggle 25 menit. Berkeliling, jangan menjelaskan materi dulu.}

\begin{frame}{Misi kalian}
\framesubtitle{Struggle, 25 menit}
\begin{columns}[T]
\begin{column}{0.31\textwidth}
\begin{Langkah}{1}{Deklarasi}
Buat \texttt{int a[5] = \{70, 85, 60, 90, 75\};} dan tampilkan \texttt{a[0]} dan \texttt{a[4]}.
\end{Langkah}
\end{column}
\begin{column}{0.31\textwidth}
\begin{Langkah}{2}{Ulangi}
Tampilkan semua elemen dengan \texttt{for}.
\end{Langkah}
\end{column}
\begin{column}{0.31\textwidth}
\begin{Langkah}{3}{Uji batas}
Tampilkan \texttt{a[5]}. Apa yang muncul?
\end{Langkah}
\end{column}
\end{columns}
\vspace{8pt}
\begin{columns}[T]
\begin{column}{0.47\textwidth}
\begin{Blok}{Boleh}
Bertanya ke teman, mengubah apa saja, sengaja membuat error.
\end{Blok}
\end{column}
\begin{column}{0.47\textwidth}
\begin{BlokAlert}{Belum dulu}
Mencari jawaban jadi di internet atau AI.
\end{BlokAlert}
\end{column}
\end{columns}
\note{Catat kesalahan yang sering muncul untuk dibahas di debrief.}
\end{frame}

\begin{frame}{Apa yang kalian temukan?}
\framesubtitle{Debrief struggle}
\begin{columns}[T]
\begin{column}{0.31\textwidth}
\begin{BlokGelap}{Pertanyaan 1}
Mengapa indeks pertama 0, bukan 1?
\end{BlokGelap}
\end{column}
\begin{column}{0.31\textwidth}
\begin{BlokGelap}{Pertanyaan 2}
Apa yang tampil untuk \texttt{a[5]}, dan apakah sama di setiap komputer?
\end{BlokGelap}
\end{column}
\begin{column}{0.31\textwidth}
\begin{BlokGelap}{Pertanyaan 3}
Bagaimana kalian menghitung jumlah semua elemen?
\end{BlokGelap}
\end{column}
\end{columns}
\vspace{10pt}
Jawaban kalian akan kita uji di bagian berikutnya.\note{Tulis jawaban di papan; buktikan atau patahkan di bagian materi.}
\end{frame}

\Bagian{3}{Banyak nilai, satu nama}{Deklarasi, indeks, perulangan, dan algoritma dasar array}
\note{Materi 40 menit. Tunjukkan setiap contoh langsung di GDB Online.}

\begin{frame}[fragile]{Deklarasi dan indeks}
\framesubtitle{Array satu dimensi}
\begin{columns}[T]
\begin{column}{0.55\textwidth}
\begin{Kode}[]{array.cpp}
int nilai[5] = {70, 85, 60, 90, 75};
cout << nilai[0] << " " << nilai[4] << endl;
nilai[2] = 65;
nilai[3] += 5;
for (int i = 0; i < 5; i++) {
    cout << nilai[i] << " ";
}
cout << endl;
\end{Kode}
\end{column}
\begin{column}{0.41\textwidth}
\only<1>{\begin{TebakDulu}
{\ITERAKickerFont\fontsize{10pt}{12pt}\selectfont\color{ITERAAlert}TEBAK DULU}\par\vspace{4pt}
Sebelum dijalankan, tulis keluarannya di kertas.
\end{TebakDulu}
}\only<2>{\KeluaranFile[]{keluaran/P09/k001.txt}
\begin{Catatan}
Array berukuran 5 punya indeks 0 sampai 4. Ukuran harus konstanta.
\end{Catatan}
}
\end{column}
\end{columns}
\note<1>{Minta mahasiswa menebak keluaran sebelum membuka slide berikutnya. Tanyakan satu atau dua tebakan yang berbeda, lalu buktikan.}
\end{frame}

\begin{frame}[fragile]{Mengisi array dari masukan}
\framesubtitle{Pola baca-proses}
\begin{columns}[T]
\begin{column}{0.55\textwidth}
\begin{Kode}[]{baca.cpp}
int n, a[100];
cin >> n;
for (int i = 0; i < n; i++) {
    cin >> a[i];
}
int total = 0;
for (int i = 0; i < n; i++) {
    total += a[i];
}
cout << "Total " << total << endl;
\end{Kode}
\end{column}
\begin{column}{0.41\textwidth}
\begin{Keluaran}[title={Masukan}]
4
10 20 30 40
\end{Keluaran}
\only<1>{\begin{TebakDulu}
{\ITERAKickerFont\fontsize{10pt}{12pt}\selectfont\color{ITERAAlert}TEBAK DULU}\par\vspace{4pt}
Sebelum dijalankan, tulis keluarannya di kertas.
\end{TebakDulu}
}\only<2>{\KeluaranFile[]{keluaran/P09/k002.txt}
\begin{Catatan}
Deklarasikan array cukup besar (misalnya 100), lalu pakai hanya n elemen pertama.
\end{Catatan}
}
\end{column}
\end{columns}
\note<1>{Minta mahasiswa menebak keluaran sebelum membuka slide berikutnya. Tanyakan satu atau dua tebakan yang berbeda, lalu buktikan.}
\end{frame}

\begin{frame}[fragile]{Mencari maksimum dan posisinya}
\framesubtitle{Algoritma dasar}
\begin{columns}[T]
\begin{column}{0.60\textwidth}
\begin{Kode}[listing options={style=ITERACodeKecil}]{maks.cpp}
int a[6] = {70, 85, 60, 90, 75, 80};
int iMaks = 0;
for (int i = 1; i < 6; i++) {
    if (a[i] > a[iMaks]) {
        iMaks = i;
    }
}
cout << a[iMaks] << " di indeks " << iMaks << endl;
\end{Kode}
\end{column}
\begin{column}{0.36\textwidth}
\only<1>{\begin{TebakDulu}
{\ITERAKickerFont\fontsize{10pt}{12pt}\selectfont\color{ITERAAlert}TEBAK DULU}\par\vspace{4pt}
Sebelum dijalankan, tulis keluarannya di kertas.
\end{TebakDulu}
}\only<2>{\KeluaranFile[listing options={style=ITERAPlainKecil}]{keluaran/P09/k003.txt}
\begin{Catatan}
Simpan indeks, bukan nilainya. Dari indeks, nilai selalu bisa diambil kembali.
\end{Catatan}
}
\end{column}
\end{columns}
\note<1>{Minta mahasiswa menebak keluaran sebelum membuka slide berikutnya. Tanyakan satu atau dua tebakan yang berbeda, lalu buktikan.}
\end{frame}

\begin{frame}[fragile]{Dua putaran: rata-rata lalu bandingkan}
\framesubtitle{Mengapa perlu array}
\begin{columns}[T]
\begin{column}{0.55\textwidth}
\begin{Kode}[]{diatas.cpp}
int a[6] = {70, 85, 60, 90, 75, 80};
int total = 0;
for (int i = 0; i < 6; i++) total += a[i];
double rata = total / 6.0;
int diAtas = 0;
for (int i = 0; i < 6; i++) {
    if (a[i] > rata) diAtas++;
}
cout << rata << " " << diAtas << endl;
\end{Kode}
\end{column}
\begin{column}{0.41\textwidth}
\only<1>{\begin{TebakDulu}
{\ITERAKickerFont\fontsize{10pt}{12pt}\selectfont\color{ITERAAlert}TEBAK DULU}\par\vspace{4pt}
Sebelum dijalankan, tulis keluarannya di kertas.
\end{TebakDulu}
}\only<2>{\KeluaranFile[]{keluaran/P09/k004.txt}
}
\end{column}
\end{columns}
\note<1>{Minta mahasiswa menebak keluaran sebelum membuka slide berikutnya. Tanyakan satu atau dua tebakan yang berbeda, lalu buktikan.}
\end{frame}

\begin{frame}[fragile]{Di luar batas}
\framesubtitle{Kesalahan paling berbahaya}
\begin{columns}[T]
\begin{column}{0.55\textwidth}
\begin{Kode}[]{batas.cpp}
int a[3] = {1, 2, 3};
int jumlah = 0;
for (int i = 0; i < 3; i++) {
    jumlah += a[i];
}
cout << jumlah << endl;
// for (int i = 0; i <= 3; i++) membaca a[3]:
// di luar array, hasilnya tidak bisa ditebak
\end{Kode}
\end{column}
\begin{column}{0.41\textwidth}
\only<1>{\begin{TebakDulu}
{\ITERAKickerFont\fontsize{10pt}{12pt}\selectfont\color{ITERAAlert}TEBAK DULU}\par\vspace{4pt}
Sebelum dijalankan, tulis keluarannya di kertas.
\end{TebakDulu}
}\only<2>{\KeluaranFile[]{keluaran/P09/k005.txt}
\begin{Catatan}
C++ tidak memeriksa batas indeks. Program bisa tetap jalan dengan hasil acak, atau berhenti mendadak.
\end{Catatan}
}
\end{column}
\end{columns}
\note<1>{Minta mahasiswa menebak keluaran sebelum membuka slide berikutnya. Tanyakan satu atau dua tebakan yang berbeda, lalu buktikan.}
\end{frame}

\begin{frame}{Menggeser dan membalik}
\framesubtitle{Menelusuri isi array}
Rotasi kanan \{10, 20, 30, 40, 50\}: simpan elemen terakhir, geser dari belakang, isi depan.
\vspace{8pt}

\small
\begin{tabular}{@{}>{\raggedright\arraybackslash}p{172pt}>{\raggedright\arraybackslash}p{364pt}>{\raggedright\arraybackslash}p{326pt}@{}}
\kepala{Langkah} & \kepala{Isi array} & \kepala{Keterangan}\\
awal & 10 20 30 40 50 & simpan = 50\\
\barisgenap i = 4 & 10 20 30 40 40 & a[4] = a[3]\\
i = 3 & 10 20 30 30 40 & a[3] = a[2]\\
\barisgenap i = 1 & 10 10 20 30 40 & setelah i = 2 dan i = 1\\
akhir & 50 10 20 30 40 & a[0] = simpan\\
\garisdasar
\end{tabular}
\end{frame}

\begin{frame}[fragile]{Tebak keluarannya}
\framesubtitle{Latihan menjelaskan, 3 menit}
\begin{columns}[T]
\begin{column}{0.56\textwidth}
\begin{Kode}[]{tebak.cpp}
int a[5] = {3, 1, 4, 1, 5};
for (int i = 1; i < 5; i++) {
    a[i] = a[i] + a[i - 1];
}
for (int i = 0; i < 5; i++) {
    cout << a[i] << " ";
}
cout << endl;
\end{Kode}
\end{column}
\begin{column}{0.40\textwidth}
\begin{Catatan}
Tulis keluarannya di kertas, karakter demi karakter. Jangan dijalankan dulu.
\end{Catatan}
\end{column}
\end{columns}
\note{Contoh soal Bagian B. Beri waktu 3 menit sebelum slide jawaban.}
\end{frame}

\begin{frame}[fragile]{Tebak keluarannya: jawaban}
\framesubtitle{Latihan menjelaskan}
\begin{columns}[T]
\begin{column}{0.56\textwidth}
\begin{Kode}[]{tebak.cpp}
int a[5] = {3, 1, 4, 1, 5};
for (int i = 1; i < 5; i++) {
    a[i] = a[i] + a[i - 1];
}
for (int i = 0; i < 5; i++) {
    cout << a[i] << " ";
}
cout << endl;
\end{Kode}
\end{column}
\begin{column}{0.40\textwidth}
\begin{Keluaran}[]
3 4 8 9 14 
\end{Keluaran}
\begin{Catatan}
Setiap elemen ditambah elemen sebelumnya yang sudah diperbarui, sehingga terbentuk jumlah kumulatif: 3, 4, 8, 9, 14.
\end{Catatan}
\end{column}
\end{columns}
\note{Tanyakan jawaban yang paling sering salah, lalu telusuri bersama.}
\end{frame}

\begin{frame}{Error yang sering muncul minggu ini}
\framesubtitle{Pesan diambil langsung dari g++}
\footnotesize
\begin{tabular}{@{}>{\raggedright\arraybackslash}p{208pt}>{\raggedright\arraybackslash}p{256pt}>{\raggedright\arraybackslash}p{398pt}@{}}
\kepala{Kesalahan} & \kepala{Contoh} & \kepala{Pesan pertama dari g++}\\
Ukuran negatif & \texttt{int a[-3];} & \texttt{error: narrowing conversion of '-3' from 'int' to 'long unsigned int'}\\
\barisgenap Inisialisasi kelebihan elemen & \texttt{int a[2] = \{1, 2, 3\};} & \texttt{error: too many initializers for 'int [2]'}\\
Indeks bukan bilangan bulat & \texttt{a[1.5]} & \texttt{error: invalid types 'int [3][double]' for array subscript}\\
\barisgenap Menugaskan array ke array & \texttt{b = a;} & \texttt{error: invalid array assignment}\\
Ukuran tidak ditulis & \texttt{int a[];} & \texttt{error: storage size of 'a' isn't known}\\
\garisdasar
\end{tabular}
\vspace{10pt}

{\small Pesan di tabel ini dihasilkan dengan g++ dan opsi -Wall, sama seperti di cek.sh.}
\note{Minta mahasiswa memotret slide ini.}
\end{frame}

\Bagian{4}{Hands-on bertingkat}{45 menit, dari dasar sampai tantangan}
\note{Mahasiswa membuka folder 02 Hands-on/P9 - Array 1D - Hands-on.}

\begin{frame}{Peta hands-on}
\framesubtitle{Folder 02 Hands-on/P9 - Array 1D - Hands-on}
\small
\begin{tabular}{@{}>{\raggedright\arraybackslash}p{36pt}>{\raggedright\arraybackslash}p{267pt}>{\raggedright\arraybackslash}p{110pt}>{\raggedright\arraybackslash}p{83pt}>{\raggedright\arraybackslash}p{341pt}@{}}
\kepala{No} & \kepala{File} & \kepala{Tingkat} & \kepala{Waktu} & \kepala{Yang dilatih}\\
1 & \texttt{handson1-nilai.cpp} & Dasar & 10' & mengisi dan menampilkan, urutan terbalik\\
\barisgenap 2 & \texttt{handson2-statistik.cpp} & Menengah & 15' & rata-rata, maksimum, dua putaran\\
3 & \texttt{handson3-geser.cpp} & Tantangan & 20' & telusuri isi array, perbaiki tiga kesalahan indeks\\
\barisgenap + & \texttt{bonus-frekuensi.cpp} & Pengayaan & sisa & array sebagai pencacah\\
\garisdasar
\end{tabular}
\vspace{10pt}

Kerjakan berurutan. Cocokkan keluaran dengan folder \texttt{expected/}; latihan yang membaca masukan memakai data di folder \texttt{input/}.\note{Saat berkeliling, minta mahasiswa yang mengerjakan latihan tantangan menjelaskan blok PENJELASAN mereka.}
\end{frame}

\begin{frame}{Cara mengecek dan aturannya}
\framesubtitle{Hands-on}
\begin{columns}[T]
\begin{column}{0.47\textwidth}
\begin{Blok}{Di GDB Online}
Salin isi file latihan, lengkapi TODO, klik Run. Bila program membaca masukan, ketik isi file di \texttt{input/}. Bandingkan dengan \texttt{expected/}.
\end{Blok}
\end{column}
\begin{column}{0.47\textwidth}
\begin{BlokGelap}{Di komputer sendiri}
Jalankan \texttt{bash cek.sh}. Skrip mengompilasi dengan \texttt{-Wall -Werror}, memberi masukan dari \texttt{input/}, lalu mencocokkan keluaran.
\end{BlokGelap}
\end{column}
\end{columns}
\vspace{6pt}
\begin{Catatan}
AI boleh untuk memahami pesan error, tidak untuk meminta solusi utuh. Exit Ticket dikerjakan tanpa bantuan apa pun.
\end{Catatan}
\note{Hands-on tidak dinilai langsung; yang dinilai Exit Ticket dan tugas.}
\end{frame}

\Bagian{5}{Exit Ticket dan tugas}{25 menit terakhir, lalu pekerjaan rumah}
\note{Mulai Exit Ticket tepat waktu.}

\begin{frame}{Exit Ticket 9}
\framesubtitle{Individual, 25 menit, tanpa AI dan internet}
\small
\begin{tabular}{@{}>{\raggedright\arraybackslash}p{60pt}>{\raggedright\arraybackslash}p{560pt}>{\raggedright\arraybackslash}p{80pt}>{\raggedright\arraybackslash}p{150pt}@{}}
\kepala{Soal} & \kepala{Isi} & \kepala{Poin} & \kepala{CPMK}\\
A1 & Tulis program yang membaca n nilai ke array dan menampilkan nilai minimum beserta indeksnya & 25 & CPMK0607\\
\barisgenap A2 & Tulis program yang menghitung banyak elemen genap dan ganjil & 25 & CPMK0607\\
B1 & Telusuri isi array setelah sebuah perulangan yang mengubah elemennya & 20 & CPMK0608\\
\barisgenap B2 & Jelaskan mengapa sebuah perulangan menyebabkan akses di luar batas & 20 & CPMK0608\\
B3 & Jelaskan mengapa sebuah masalah butuh array, bukan variabel tunggal & 10 & CPMK0608\\
\garisdasar
\end{tabular}
\vspace{10pt}

{\small Soal A dinilai dari keluaran yang tepat sama. Soal B dinilai dari ketepatan dan alasan. Pelanggaran aturan membuat nilai Exit Ticket ini 0.}
\note{Soal lengkap dibagikan terpisah dan tidak disimpan di repo publik.}
\end{frame}

\begin{frame}{Tugas 9: Sistem Nilai Mahasiswa}
\framesubtitle{Dikumpulkan sebelum Pertemuan 10}
\begin{columns}[T]
\begin{column}{0.47\textwidth}
\begin{Blok}{Bagian A, program, 50 poin}
Buat program dengan ketentuan berikut. Gunakan \texttt{const} untuk ukuran array. Ketentuannya di slide berikut.
\end{Blok}
\end{column}
\begin{column}{0.47\textwidth}
\begin{BlokGelap}{Bagian B, penjelasan, 50 poin}
Komentar maksimal 150 kata dan gambar isi kedua array setelah membaca tiga data contoh. Jelaskan cara menjaga indeks tetap di dalam batas, dan mengapa dua array harus diproses dengan indeks yang sama.
\end{BlokGelap}
\end{column}
\end{columns}
\vspace{8pt}
{\small Data di tugas adalah data kalian sendiri. Rubrik lengkap ada di Modul Belajar 9.}
\note{Ingatkan bahwa penjelasan disalin dari AI mudah dikenali karena tidak menyebut pengalaman mereka sendiri.}
\end{frame}

\begin{frame}{Ketentuan Tugas 9}
\framesubtitle{Sistem Nilai Mahasiswa}
\small
\begin{tabular}{@{}>{\raggedright\arraybackslash}p{87pt}>{\raggedright\arraybackslash}p{788pt}@{}}
\kepala{No} & \kepala{Ketentuan Bagian A}\\
1 & Baca jumlah mahasiswa (maksimal 50)\\
\barisgenap 2 & Baca NIM dan nilai (0 sampai 100) setiap mahasiswa ke dua array sejajar, dengan validasi\\
3 & Tampilkan semua data, nilai tertinggi dan terendah, serta rata-rata kelas\\
\barisgenap 4 & Tampilkan jumlah LULUS (nilai 60 ke atas), TIDAK LULUS, dan persentase kelulusan\\
5 & Fitur pencarian: baca sebuah NIM, tampilkan nilainya atau pesan tidak ditemukan\\
\barisgenap Bonus & Urutkan nilai dari tertinggi ke terendah (materi sorting di Pertemuan 12).\\
\garisdasar
\end{tabular}
\note{Bacakan ketentuan satu per satu dan tanyakan apakah ada yang belum jelas.}
\end{frame}

\begin{frame}{Rubrik Tugas 9}
\framesubtitle{Empat level, sama di semua instrumen}
\footnotesize
\begin{tabular}{@{}>{\raggedright\arraybackslash}p{166pt}>{\raggedright\arraybackslash}p{44pt}>{\raggedright\arraybackslash}p{166pt}>{\raggedright\arraybackslash}p{156pt}>{\raggedright\arraybackslash}p{146pt}>{\raggedright\arraybackslash}p{146pt}@{}}
\kepala{Kriteria} & \kepala{Poin} & \kepala{Sangat baik 85--100} & \kepala{Baik 70--84} & \kepala{Cukup 55--69} & \kepala{Kurang <55}\\
A1 Program berjalan & 20 & tanpa error dan peringatan & ada peringatan & perlu satu perbaikan & tidak terkompilasi\\
\barisgenap A2 Kelengkapan fitur & 15 & semua statistik dan pencarian & satu fitur kurang & dua kurang & tiga atau lebih kurang\\
A3 Validasi dan ketepatan & 15 & validasi tepat, hasil benar untuk n = 1 dan n = 10 & satu salah & dua salah & sebagian besar salah\\
\barisgenap B1 Isi array dan batas indeks & 30 & gambar tepat, alasan jelas & satu langkah keliru & gambar tanpa alasan & tidak ada atau disalin\\
B2 Cerita error dan pengujian & 20 & pesan, sebab, perbaikan, uji & pesan tidak dikutip & hanya perbaikan & tidak ada\\
\garisdasar
\end{tabular}
\vspace{8pt}

{\small Nilai kriteria = poin × skor level. Bagian A total 50, Bagian B total 50.}
\note{Rubrik ini sama strukturnya setiap minggu; yang berubah hanya isi kriteria A2, A3, dan B1.}
\end{frame}

\begin{frame}{Sebelum Pertemuan 10}
\framesubtitle{Persiapan}
\begin{columns}[T]
\begin{column}{0.31\textwidth}
\begin{Langkah}{1}{Selesaikan}
Hands-on yang belum lulus \texttt{cek.sh}, terutama latihan tantangan.
\end{Langkah}
\end{column}
\begin{column}{0.31\textwidth}
\begin{Langkah}{2}{Kumpulkan}
Tugas 9 sebelum Pertemuan 10 dimulai.
\end{Langkah}
\end{column}
\begin{column}{0.31\textwidth}
\begin{Langkah}{3}{Baca}
Modul Belajar 9 untuk mengulang, lalu bagian awal modul berikutnya.
\end{Langkah}
\end{column}
\end{columns}
\vspace{10pt}
Pertemuan 10 membahas array dua dimensi dan string, sejalan dengan Algoritma Pertemuan 10.\note{Tanyakan satu hal yang paling membingungkan hari ini untuk dibuka di review minggu depan.}
\end{frame}

\SlidePenutup{Terima kasih}{Sampai jumpa di Pertemuan 10: array dua dimensi dan string}
\note{Slide penutup.}
\end{document}
"
\documentclass[t]{beamer}
\usetheme{ITERA}
% Mode catatan pembicara: aktif bila \catatan didefinisikan sebelum \input.
\ifdefined\catatan
  \setbeameroption{show only notes}
  \setbeamertemplate{note page}{\vspace*{20pt}\hspace*{4pt}\begin{minipage}{900pt}
    {\ITERAKickerFont\fontsize{11pt}{14pt}\selectfont\color{ITERAGoldText}CATATAN PEMBICARA \ \ |\ \ SLIDE \insertframenumber}\par\vspace{4pt}
    {\ITERADisplay\bfseries\fontsize{22pt}{28pt}\selectfont\color{ITERAStructure}\insertframetitle}\par\vspace{12pt}
    \fontsize{15pt}{23pt}\selectfont\insertnote\end{minipage}}
\fi
\tikzset{fase/.style={draw=none,fill=#1,rounded corners=3pt,minimum width=158pt,minimum height=118pt,text width=146pt,align=center}}

\renewcommand\ITERAmeeting{Pertemuan 9}
\begin{document}
\SlideJudul{IF25-11002 PRAKTIKUM PEMROGRAMAN}{Pertemuan 9\\Array Satu Dimensi}{Menyimpan banyak nilai sejenis, mengaksesnya dengan indeks, dan memprosesnya dengan perulangan}{Tim Pengajar}{Tahun 2026. Sejajar dengan Pertemuan 9 Algoritma Pemrograman: Array 1D.}
\note{Buka dengan menanyakan satu hal dari pertemuan lalu yang masih membingungkan.}

\begin{frame}{Dua mata kuliah, satu minggu}
\framesubtitle{Sebelum mulai}
Topik minggu ini sama di dua mata kuliah, dengan peran yang berbeda.
\vspace{10pt}

\begin{columns}[T]
\begin{column}{0.47\textwidth}
\begin{Blok}{Algoritma: Array 1D}
Konsep larik, indeks, dan algoritma dasar pada larik: penjumlahan, maksimum, dan pembalikan.
\end{Blok}
\end{column}
\begin{column}{0.47\textwidth}
\begin{BlokGelap}{Praktikum: Array 1D}
Mendeklarasikan array C++, mengisi dan membacanya dengan \texttt{for}, serta menghindari akses di luar batas.
\end{BlokGelap}
\end{column}
\end{columns}
\note{Tanyakan apa yang sudah dibahas di Algoritma minggu ini. Gunakan jawabannya sebagai jembatan ke praktikum.}
\end{frame}

\begin{frame}{Saat pulang nanti, kalian bisa}
\framesubtitle{Target hari ini}
\begin{columns}[T]
\begin{column}{0.47\textwidth}
\begin{Langkah}{1}{Menerapkan}
Menulis program yang menyimpan data sejenis di array satu dimensi dan memprosesnya dengan perulangan: mengisi, menampilkan, menjumlahkan, mencari maksimum, dan menggeser

{\footnotesize\color{ITERAInkMuted}Sub-CPMK 9.1, CPMK0607}
\end{Langkah}
\end{column}
\begin{column}{0.47\textwidth}
\begin{Langkah}{2}{Menjelaskan}
Menelusuri isi array sesudah setiap langkah pemrosesan dan menjelaskan kesalahan indeks, termasuk akses di luar batas

{\footnotesize\color{ITERAInkMuted}Sub-CPMK 9.2, CPMK0608}
\end{Langkah}
\end{column}
\end{columns}
\vspace{8pt}
\begin{Catatan}
Dua kemampuan ini dinilai sama berat di Exit Ticket, tugas, UTS, dan UAS.
\end{Catatan}
\note{Bacakan target dengan pelan. Kata kuncinya menerapkan dan menjelaskan.}
\end{frame}

\begin{frame}{Ingat minggu lalu}
\framesubtitle{Review, 5 menit}
\begin{columns}[T]
\begin{column}{0.31\textwidth}
\begin{BlokGelap}{Pertanyaan 1}
Apa hasil \texttt{int n = 4096} setelah dijumlah digitnya?
\end{BlokGelap}
\end{column}
\begin{column}{0.31\textwidth}
\begin{BlokGelap}{Pertanyaan 2}
Berapa kali badan \texttt{for i 1..3, for j 1..i} dijalankan?
\end{BlokGelap}
\end{column}
\begin{column}{0.31\textwidth}
\begin{BlokGelap}{Pertanyaan 3}
Apa beda soal A dan soal B di UTS?
\end{BlokGelap}
\end{column}
\end{columns}
\vspace{8pt}
Jawab di kertas dulu, lalu cocokkan dengan teman sebelah.\note{Pertemuan pertama sesudah UTS. Bahas singkat kesalahan umum UTS sebelum masuk materi.}
\end{frame}

\Bagian{1}{Masalah pembuka}{Siapa yang di atas rata-rata?}
\note{Masalah pembuka 10 menit.}

\begin{frame}[fragile]{Siapa yang di atas rata-rata?}
\framesubtitle{Masalah minggu ini}
\begin{columns}[T]
\begin{column}{0.55\textwidth}
{\small Dosen ingin tahu berapa mahasiswa yang nilainya di atas rata-rata kelas.

\vspace{6pt}
Dengan pola sentinel minggu 5, rata-rata baru diketahui setelah semua nilai dibaca. Tetapi saat itu nilai-nilai sebelumnya sudah hilang, karena setiap \texttt{cin} menimpa variabel yang sama.\par}
\vspace{6pt}
\begin{Catatan}
\textbf{Diskusi 2 menit.} Apa yang perlu dilakukan supaya semua nilai masih ada saat rata-rata sudah diketahui? Berapa variabel yang dibutuhkan untuk 40 mahasiswa?
\end{Catatan}
\end{column}
\begin{column}{0.41\textwidth}
\begin{Keluaran}[]
Nilai: 70 85 60 90 75 80
Rata-rata   : 76.67
Di atas rata: 3
\end{Keluaran}
\end{column}
\end{columns}
\note{Tampung beberapa jawaban tanpa menilai benar salah.}
\end{frame}

\begin{frame}{Dari langkah ke instruksi}
\framesubtitle{Sambungan dengan Algoritma}
\begin{columns}[T]
\begin{column}{0.31\textwidth}
\begin{Langkah}{1}{Simpan semua}
Semua nilai perlu disimpan, bukan hanya totalnya.
\end{Langkah}
\end{column}
\begin{column}{0.31\textwidth}
\begin{Langkah}{2}{Proses dua kali}
Putaran pertama menghitung rata-rata, putaran kedua membandingkan.
\end{Langkah}
\end{column}
\begin{column}{0.31\textwidth}
\begin{Langkah}{3}{Terjemahkan}
Satu array, dua \texttt{for} dengan indeks 0 sampai n − 1.
\end{Langkah}
\end{column}
\end{columns}
\vspace{10pt}
Langkah 1 dan 2 dilatih di Algoritma. Langkah 3 kita kerjakan hari ini dengan C++.\note{Hubungkan eksplisit ke Algoritma minggu ini.}
\end{frame}

\Bagian{2}{Coba dulu, baru dijelaskan}{25 menit eksperimen di GDB Online}
\note{Struggle 25 menit. Berkeliling, jangan menjelaskan materi dulu.}

\begin{frame}{Misi kalian}
\framesubtitle{Struggle, 25 menit}
\begin{columns}[T]
\begin{column}{0.31\textwidth}
\begin{Langkah}{1}{Deklarasi}
Buat \texttt{int a[5] = \{70, 85, 60, 90, 75\};} dan tampilkan \texttt{a[0]} dan \texttt{a[4]}.
\end{Langkah}
\end{column}
\begin{column}{0.31\textwidth}
\begin{Langkah}{2}{Ulangi}
Tampilkan semua elemen dengan \texttt{for}.
\end{Langkah}
\end{column}
\begin{column}{0.31\textwidth}
\begin{Langkah}{3}{Uji batas}
Tampilkan \texttt{a[5]}. Apa yang muncul?
\end{Langkah}
\end{column}
\end{columns}
\vspace{8pt}
\begin{columns}[T]
\begin{column}{0.47\textwidth}
\begin{Blok}{Boleh}
Bertanya ke teman, mengubah apa saja, sengaja membuat error.
\end{Blok}
\end{column}
\begin{column}{0.47\textwidth}
\begin{BlokAlert}{Belum dulu}
Mencari jawaban jadi di internet atau AI.
\end{BlokAlert}
\end{column}
\end{columns}
\note{Catat kesalahan yang sering muncul untuk dibahas di debrief.}
\end{frame}

\begin{frame}{Apa yang kalian temukan?}
\framesubtitle{Debrief struggle}
\begin{columns}[T]
\begin{column}{0.31\textwidth}
\begin{BlokGelap}{Pertanyaan 1}
Mengapa indeks pertama 0, bukan 1?
\end{BlokGelap}
\end{column}
\begin{column}{0.31\textwidth}
\begin{BlokGelap}{Pertanyaan 2}
Apa yang tampil untuk \texttt{a[5]}, dan apakah sama di setiap komputer?
\end{BlokGelap}
\end{column}
\begin{column}{0.31\textwidth}
\begin{BlokGelap}{Pertanyaan 3}
Bagaimana kalian menghitung jumlah semua elemen?
\end{BlokGelap}
\end{column}
\end{columns}
\vspace{10pt}
Jawaban kalian akan kita uji di bagian berikutnya.\note{Tulis jawaban di papan; buktikan atau patahkan di bagian materi.}
\end{frame}

\Bagian{3}{Banyak nilai, satu nama}{Deklarasi, indeks, perulangan, dan algoritma dasar array}
\note{Materi 40 menit. Tunjukkan setiap contoh langsung di GDB Online.}

\begin{frame}[fragile]{Deklarasi dan indeks}
\framesubtitle{Array satu dimensi}
\begin{columns}[T]
\begin{column}{0.55\textwidth}
\begin{Kode}[]{array.cpp}
int nilai[5] = {70, 85, 60, 90, 75};
cout << nilai[0] << " " << nilai[4] << endl;
nilai[2] = 65;
nilai[3] += 5;
for (int i = 0; i < 5; i++) {
    cout << nilai[i] << " ";
}
cout << endl;
\end{Kode}
\end{column}
\begin{column}{0.41\textwidth}
\only<1>{\begin{TebakDulu}
{\ITERAKickerFont\fontsize{10pt}{12pt}\selectfont\color{ITERAAlert}TEBAK DULU}\par\vspace{4pt}
Sebelum dijalankan, tulis keluarannya di kertas.
\end{TebakDulu}
}\only<2>{\KeluaranFile[]{keluaran/P09/k001.txt}
\begin{Catatan}
Array berukuran 5 punya indeks 0 sampai 4. Ukuran harus konstanta.
\end{Catatan}
}
\end{column}
\end{columns}
\note<1>{Minta mahasiswa menebak keluaran sebelum membuka slide berikutnya. Tanyakan satu atau dua tebakan yang berbeda, lalu buktikan.}
\end{frame}

\begin{frame}[fragile]{Mengisi array dari masukan}
\framesubtitle{Pola baca-proses}
\begin{columns}[T]
\begin{column}{0.55\textwidth}
\begin{Kode}[]{baca.cpp}
int n, a[100];
cin >> n;
for (int i = 0; i < n; i++) {
    cin >> a[i];
}
int total = 0;
for (int i = 0; i < n; i++) {
    total += a[i];
}
cout << "Total " << total << endl;
\end{Kode}
\end{column}
\begin{column}{0.41\textwidth}
\begin{Keluaran}[title={Masukan}]
4
10 20 30 40
\end{Keluaran}
\only<1>{\begin{TebakDulu}
{\ITERAKickerFont\fontsize{10pt}{12pt}\selectfont\color{ITERAAlert}TEBAK DULU}\par\vspace{4pt}
Sebelum dijalankan, tulis keluarannya di kertas.
\end{TebakDulu}
}\only<2>{\KeluaranFile[]{keluaran/P09/k002.txt}
\begin{Catatan}
Deklarasikan array cukup besar (misalnya 100), lalu pakai hanya n elemen pertama.
\end{Catatan}
}
\end{column}
\end{columns}
\note<1>{Minta mahasiswa menebak keluaran sebelum membuka slide berikutnya. Tanyakan satu atau dua tebakan yang berbeda, lalu buktikan.}
\end{frame}

\begin{frame}[fragile]{Mencari maksimum dan posisinya}
\framesubtitle{Algoritma dasar}
\begin{columns}[T]
\begin{column}{0.60\textwidth}
\begin{Kode}[listing options={style=ITERACodeKecil}]{maks.cpp}
int a[6] = {70, 85, 60, 90, 75, 80};
int iMaks = 0;
for (int i = 1; i < 6; i++) {
    if (a[i] > a[iMaks]) {
        iMaks = i;
    }
}
cout << a[iMaks] << " di indeks " << iMaks << endl;
\end{Kode}
\end{column}
\begin{column}{0.36\textwidth}
\only<1>{\begin{TebakDulu}
{\ITERAKickerFont\fontsize{10pt}{12pt}\selectfont\color{ITERAAlert}TEBAK DULU}\par\vspace{4pt}
Sebelum dijalankan, tulis keluarannya di kertas.
\end{TebakDulu}
}\only<2>{\KeluaranFile[listing options={style=ITERAPlainKecil}]{keluaran/P09/k003.txt}
\begin{Catatan}
Simpan indeks, bukan nilainya. Dari indeks, nilai selalu bisa diambil kembali.
\end{Catatan}
}
\end{column}
\end{columns}
\note<1>{Minta mahasiswa menebak keluaran sebelum membuka slide berikutnya. Tanyakan satu atau dua tebakan yang berbeda, lalu buktikan.}
\end{frame}

\begin{frame}[fragile]{Dua putaran: rata-rata lalu bandingkan}
\framesubtitle{Mengapa perlu array}
\begin{columns}[T]
\begin{column}{0.55\textwidth}
\begin{Kode}[]{diatas.cpp}
int a[6] = {70, 85, 60, 90, 75, 80};
int total = 0;
for (int i = 0; i < 6; i++) total += a[i];
double rata = total / 6.0;
int diAtas = 0;
for (int i = 0; i < 6; i++) {
    if (a[i] > rata) diAtas++;
}
cout << rata << " " << diAtas << endl;
\end{Kode}
\end{column}
\begin{column}{0.41\textwidth}
\only<1>{\begin{TebakDulu}
{\ITERAKickerFont\fontsize{10pt}{12pt}\selectfont\color{ITERAAlert}TEBAK DULU}\par\vspace{4pt}
Sebelum dijalankan, tulis keluarannya di kertas.
\end{TebakDulu}
}\only<2>{\KeluaranFile[]{keluaran/P09/k004.txt}
}
\end{column}
\end{columns}
\note<1>{Minta mahasiswa menebak keluaran sebelum membuka slide berikutnya. Tanyakan satu atau dua tebakan yang berbeda, lalu buktikan.}
\end{frame}

\begin{frame}[fragile]{Di luar batas}
\framesubtitle{Kesalahan paling berbahaya}
\begin{columns}[T]
\begin{column}{0.55\textwidth}
\begin{Kode}[]{batas.cpp}
int a[3] = {1, 2, 3};
int jumlah = 0;
for (int i = 0; i < 3; i++) {
    jumlah += a[i];
}
cout << jumlah << endl;
// for (int i = 0; i <= 3; i++) membaca a[3]:
// di luar array, hasilnya tidak bisa ditebak
\end{Kode}
\end{column}
\begin{column}{0.41\textwidth}
\only<1>{\begin{TebakDulu}
{\ITERAKickerFont\fontsize{10pt}{12pt}\selectfont\color{ITERAAlert}TEBAK DULU}\par\vspace{4pt}
Sebelum dijalankan, tulis keluarannya di kertas.
\end{TebakDulu}
}\only<2>{\KeluaranFile[]{keluaran/P09/k005.txt}
\begin{Catatan}
C++ tidak memeriksa batas indeks. Program bisa tetap jalan dengan hasil acak, atau berhenti mendadak.
\end{Catatan}
}
\end{column}
\end{columns}
\note<1>{Minta mahasiswa menebak keluaran sebelum membuka slide berikutnya. Tanyakan satu atau dua tebakan yang berbeda, lalu buktikan.}
\end{frame}

\begin{frame}{Menggeser dan membalik}
\framesubtitle{Menelusuri isi array}
Rotasi kanan \{10, 20, 30, 40, 50\}: simpan elemen terakhir, geser dari belakang, isi depan.
\vspace{8pt}

\small
\begin{tabular}{@{}>{\raggedright\arraybackslash}p{172pt}>{\raggedright\arraybackslash}p{364pt}>{\raggedright\arraybackslash}p{326pt}@{}}
\kepala{Langkah} & \kepala{Isi array} & \kepala{Keterangan}\\
awal & 10 20 30 40 50 & simpan = 50\\
\barisgenap i = 4 & 10 20 30 40 40 & a[4] = a[3]\\
i = 3 & 10 20 30 30 40 & a[3] = a[2]\\
\barisgenap i = 1 & 10 10 20 30 40 & setelah i = 2 dan i = 1\\
akhir & 50 10 20 30 40 & a[0] = simpan\\
\garisdasar
\end{tabular}
\end{frame}

\begin{frame}[fragile]{Tebak keluarannya}
\framesubtitle{Latihan menjelaskan, 3 menit}
\begin{columns}[T]
\begin{column}{0.56\textwidth}
\begin{Kode}[]{tebak.cpp}
int a[5] = {3, 1, 4, 1, 5};
for (int i = 1; i < 5; i++) {
    a[i] = a[i] + a[i - 1];
}
for (int i = 0; i < 5; i++) {
    cout << a[i] << " ";
}
cout << endl;
\end{Kode}
\end{column}
\begin{column}{0.40\textwidth}
\begin{Catatan}
Tulis keluarannya di kertas, karakter demi karakter. Jangan dijalankan dulu.
\end{Catatan}
\end{column}
\end{columns}
\note{Contoh soal Bagian B. Beri waktu 3 menit sebelum slide jawaban.}
\end{frame}

\begin{frame}[fragile]{Tebak keluarannya: jawaban}
\framesubtitle{Latihan menjelaskan}
\begin{columns}[T]
\begin{column}{0.56\textwidth}
\begin{Kode}[]{tebak.cpp}
int a[5] = {3, 1, 4, 1, 5};
for (int i = 1; i < 5; i++) {
    a[i] = a[i] + a[i - 1];
}
for (int i = 0; i < 5; i++) {
    cout << a[i] << " ";
}
cout << endl;
\end{Kode}
\end{column}
\begin{column}{0.40\textwidth}
\begin{Keluaran}[]
3 4 8 9 14 
\end{Keluaran}
\begin{Catatan}
Setiap elemen ditambah elemen sebelumnya yang sudah diperbarui, sehingga terbentuk jumlah kumulatif: 3, 4, 8, 9, 14.
\end{Catatan}
\end{column}
\end{columns}
\note{Tanyakan jawaban yang paling sering salah, lalu telusuri bersama.}
\end{frame}

\begin{frame}{Error yang sering muncul minggu ini}
\framesubtitle{Pesan diambil langsung dari g++}
\footnotesize
\begin{tabular}{@{}>{\raggedright\arraybackslash}p{208pt}>{\raggedright\arraybackslash}p{256pt}>{\raggedright\arraybackslash}p{398pt}@{}}
\kepala{Kesalahan} & \kepala{Contoh} & \kepala{Pesan pertama dari g++}\\
Ukuran negatif & \texttt{int a[-3];} & \texttt{error: narrowing conversion of '-3' from 'int' to 'long unsigned int'}\\
\barisgenap Inisialisasi kelebihan elemen & \texttt{int a[2] = \{1, 2, 3\};} & \texttt{error: too many initializers for 'int [2]'}\\
Indeks bukan bilangan bulat & \texttt{a[1.5]} & \texttt{error: invalid types 'int [3][double]' for array subscript}\\
\barisgenap Menugaskan array ke array & \texttt{b = a;} & \texttt{error: invalid array assignment}\\
Ukuran tidak ditulis & \texttt{int a[];} & \texttt{error: storage size of 'a' isn't known}\\
\garisdasar
\end{tabular}
\vspace{10pt}

{\small Pesan di tabel ini dihasilkan dengan g++ dan opsi -Wall, sama seperti di cek.sh.}
\note{Minta mahasiswa memotret slide ini.}
\end{frame}

\Bagian{4}{Hands-on bertingkat}{45 menit, dari dasar sampai tantangan}
\note{Mahasiswa membuka folder 02 Hands-on/P9 - Array 1D - Hands-on.}

\begin{frame}{Peta hands-on}
\framesubtitle{Folder 02 Hands-on/P9 - Array 1D - Hands-on}
\small
\begin{tabular}{@{}>{\raggedright\arraybackslash}p{36pt}>{\raggedright\arraybackslash}p{267pt}>{\raggedright\arraybackslash}p{110pt}>{\raggedright\arraybackslash}p{83pt}>{\raggedright\arraybackslash}p{341pt}@{}}
\kepala{No} & \kepala{File} & \kepala{Tingkat} & \kepala{Waktu} & \kepala{Yang dilatih}\\
1 & \texttt{handson1-nilai.cpp} & Dasar & 10' & mengisi dan menampilkan, urutan terbalik\\
\barisgenap 2 & \texttt{handson2-statistik.cpp} & Menengah & 15' & rata-rata, maksimum, dua putaran\\
3 & \texttt{handson3-geser.cpp} & Tantangan & 20' & telusuri isi array, perbaiki tiga kesalahan indeks\\
\barisgenap + & \texttt{bonus-frekuensi.cpp} & Pengayaan & sisa & array sebagai pencacah\\
\garisdasar
\end{tabular}
\vspace{10pt}

Kerjakan berurutan. Cocokkan keluaran dengan folder \texttt{expected/}; latihan yang membaca masukan memakai data di folder \texttt{input/}.\note{Saat berkeliling, minta mahasiswa yang mengerjakan latihan tantangan menjelaskan blok PENJELASAN mereka.}
\end{frame}

\begin{frame}{Cara mengecek dan aturannya}
\framesubtitle{Hands-on}
\begin{columns}[T]
\begin{column}{0.47\textwidth}
\begin{Blok}{Di GDB Online}
Salin isi file latihan, lengkapi TODO, klik Run. Bila program membaca masukan, ketik isi file di \texttt{input/}. Bandingkan dengan \texttt{expected/}.
\end{Blok}
\end{column}
\begin{column}{0.47\textwidth}
\begin{BlokGelap}{Di komputer sendiri}
Jalankan \texttt{bash cek.sh}. Skrip mengompilasi dengan \texttt{-Wall -Werror}, memberi masukan dari \texttt{input/}, lalu mencocokkan keluaran.
\end{BlokGelap}
\end{column}
\end{columns}
\vspace{6pt}
\begin{Catatan}
AI boleh untuk memahami pesan error, tidak untuk meminta solusi utuh. Exit Ticket dikerjakan tanpa bantuan apa pun.
\end{Catatan}
\note{Hands-on tidak dinilai langsung; yang dinilai Exit Ticket dan tugas.}
\end{frame}

\Bagian{5}{Exit Ticket dan tugas}{25 menit terakhir, lalu pekerjaan rumah}
\note{Mulai Exit Ticket tepat waktu.}

\begin{frame}{Exit Ticket 9}
\framesubtitle{Individual, 25 menit, tanpa AI dan internet}
\small
\begin{tabular}{@{}>{\raggedright\arraybackslash}p{60pt}>{\raggedright\arraybackslash}p{560pt}>{\raggedright\arraybackslash}p{80pt}>{\raggedright\arraybackslash}p{150pt}@{}}
\kepala{Soal} & \kepala{Isi} & \kepala{Poin} & \kepala{CPMK}\\
A1 & Tulis program yang membaca n nilai ke array dan menampilkan nilai minimum beserta indeksnya & 25 & CPMK0607\\
\barisgenap A2 & Tulis program yang menghitung banyak elemen genap dan ganjil & 25 & CPMK0607\\
B1 & Telusuri isi array setelah sebuah perulangan yang mengubah elemennya & 20 & CPMK0608\\
\barisgenap B2 & Jelaskan mengapa sebuah perulangan menyebabkan akses di luar batas & 20 & CPMK0608\\
B3 & Jelaskan mengapa sebuah masalah butuh array, bukan variabel tunggal & 10 & CPMK0608\\
\garisdasar
\end{tabular}
\vspace{10pt}

{\small Soal A dinilai dari keluaran yang tepat sama. Soal B dinilai dari ketepatan dan alasan. Pelanggaran aturan membuat nilai Exit Ticket ini 0.}
\note{Soal lengkap dibagikan terpisah dan tidak disimpan di repo publik.}
\end{frame}

\begin{frame}{Tugas 9: Sistem Nilai Mahasiswa}
\framesubtitle{Dikumpulkan sebelum Pertemuan 10}
\begin{columns}[T]
\begin{column}{0.47\textwidth}
\begin{Blok}{Bagian A, program, 50 poin}
Buat program dengan ketentuan berikut. Gunakan \texttt{const} untuk ukuran array. Ketentuannya di slide berikut.
\end{Blok}
\end{column}
\begin{column}{0.47\textwidth}
\begin{BlokGelap}{Bagian B, penjelasan, 50 poin}
Komentar maksimal 150 kata dan gambar isi kedua array setelah membaca tiga data contoh. Jelaskan cara menjaga indeks tetap di dalam batas, dan mengapa dua array harus diproses dengan indeks yang sama.
\end{BlokGelap}
\end{column}
\end{columns}
\vspace{8pt}
{\small Data di tugas adalah data kalian sendiri. Rubrik lengkap ada di Modul Belajar 9.}
\note{Ingatkan bahwa penjelasan disalin dari AI mudah dikenali karena tidak menyebut pengalaman mereka sendiri.}
\end{frame}

\begin{frame}{Ketentuan Tugas 9}
\framesubtitle{Sistem Nilai Mahasiswa}
\small
\begin{tabular}{@{}>{\raggedright\arraybackslash}p{87pt}>{\raggedright\arraybackslash}p{788pt}@{}}
\kepala{No} & \kepala{Ketentuan Bagian A}\\
1 & Baca jumlah mahasiswa (maksimal 50)\\
\barisgenap 2 & Baca NIM dan nilai (0 sampai 100) setiap mahasiswa ke dua array sejajar, dengan validasi\\
3 & Tampilkan semua data, nilai tertinggi dan terendah, serta rata-rata kelas\\
\barisgenap 4 & Tampilkan jumlah LULUS (nilai 60 ke atas), TIDAK LULUS, dan persentase kelulusan\\
5 & Fitur pencarian: baca sebuah NIM, tampilkan nilainya atau pesan tidak ditemukan\\
\barisgenap Bonus & Urutkan nilai dari tertinggi ke terendah (materi sorting di Pertemuan 12).\\
\garisdasar
\end{tabular}
\note{Bacakan ketentuan satu per satu dan tanyakan apakah ada yang belum jelas.}
\end{frame}

\begin{frame}{Rubrik Tugas 9}
\framesubtitle{Empat level, sama di semua instrumen}
\footnotesize
\begin{tabular}{@{}>{\raggedright\arraybackslash}p{166pt}>{\raggedright\arraybackslash}p{44pt}>{\raggedright\arraybackslash}p{166pt}>{\raggedright\arraybackslash}p{156pt}>{\raggedright\arraybackslash}p{146pt}>{\raggedright\arraybackslash}p{146pt}@{}}
\kepala{Kriteria} & \kepala{Poin} & \kepala{Sangat baik 85--100} & \kepala{Baik 70--84} & \kepala{Cukup 55--69} & \kepala{Kurang <55}\\
A1 Program berjalan & 20 & tanpa error dan peringatan & ada peringatan & perlu satu perbaikan & tidak terkompilasi\\
\barisgenap A2 Kelengkapan fitur & 15 & semua statistik dan pencarian & satu fitur kurang & dua kurang & tiga atau lebih kurang\\
A3 Validasi dan ketepatan & 15 & validasi tepat, hasil benar untuk n = 1 dan n = 10 & satu salah & dua salah & sebagian besar salah\\
\barisgenap B1 Isi array dan batas indeks & 30 & gambar tepat, alasan jelas & satu langkah keliru & gambar tanpa alasan & tidak ada atau disalin\\
B2 Cerita error dan pengujian & 20 & pesan, sebab, perbaikan, uji & pesan tidak dikutip & hanya perbaikan & tidak ada\\
\garisdasar
\end{tabular}
\vspace{8pt}

{\small Nilai kriteria = poin × skor level. Bagian A total 50, Bagian B total 50.}
\note{Rubrik ini sama strukturnya setiap minggu; yang berubah hanya isi kriteria A2, A3, dan B1.}
\end{frame}

\begin{frame}{Sebelum Pertemuan 10}
\framesubtitle{Persiapan}
\begin{columns}[T]
\begin{column}{0.31\textwidth}
\begin{Langkah}{1}{Selesaikan}
Hands-on yang belum lulus \texttt{cek.sh}, terutama latihan tantangan.
\end{Langkah}
\end{column}
\begin{column}{0.31\textwidth}
\begin{Langkah}{2}{Kumpulkan}
Tugas 9 sebelum Pertemuan 10 dimulai.
\end{Langkah}
\end{column}
\begin{column}{0.31\textwidth}
\begin{Langkah}{3}{Baca}
Modul Belajar 9 untuk mengulang, lalu bagian awal modul berikutnya.
\end{Langkah}
\end{column}
\end{columns}
\vspace{10pt}
Pertemuan 10 membahas array dua dimensi dan string, sejalan dengan Algoritma Pertemuan 10.\note{Tanyakan satu hal yang paling membingungkan hari ini untuk dibuka di review minggu depan.}
\end{frame}

\SlidePenutup{Terima kasih}{Sampai jumpa di Pertemuan 10: array dua dimensi dan string}
\note{Slide penutup.}
\end{document}
"
\documentclass[t]{beamer}
\usetheme{ITERA}
% Mode catatan pembicara: aktif bila \catatan didefinisikan sebelum \input.
\ifdefined\catatan
  \setbeameroption{show only notes}
  \setbeamertemplate{note page}{\vspace*{20pt}\hspace*{4pt}\begin{minipage}{900pt}
    {\ITERAKickerFont\fontsize{11pt}{14pt}\selectfont\color{ITERAGoldText}CATATAN PEMBICARA \ \ |\ \ SLIDE \insertframenumber}\par\vspace{4pt}
    {\ITERADisplay\bfseries\fontsize{22pt}{28pt}\selectfont\color{ITERAStructure}\insertframetitle}\par\vspace{12pt}
    \fontsize{15pt}{23pt}\selectfont\insertnote\end{minipage}}
\fi
\tikzset{fase/.style={draw=none,fill=#1,rounded corners=3pt,minimum width=158pt,minimum height=118pt,text width=146pt,align=center}}

\renewcommand\ITERAmeeting{Pertemuan 9}
\begin{document}
\SlideJudul{IF25-11002 PRAKTIKUM PEMROGRAMAN}{Pertemuan 9\\Array Satu Dimensi}{Menyimpan banyak nilai sejenis, mengaksesnya dengan indeks, dan memprosesnya dengan perulangan}{Tim Pengajar}{Tahun 2026. Sejajar dengan Pertemuan 9 Algoritma Pemrograman: Array 1D.}
\note{Buka dengan menanyakan satu hal dari pertemuan lalu yang masih membingungkan.}

\begin{frame}{Dua mata kuliah, satu minggu}
\framesubtitle{Sebelum mulai}
Topik minggu ini sama di dua mata kuliah, dengan peran yang berbeda.
\vspace{10pt}

\begin{columns}[T]
\begin{column}{0.47\textwidth}
\begin{Blok}{Algoritma: Array 1D}
Konsep larik, indeks, dan algoritma dasar pada larik: penjumlahan, maksimum, dan pembalikan.
\end{Blok}
\end{column}
\begin{column}{0.47\textwidth}
\begin{BlokGelap}{Praktikum: Array 1D}
Mendeklarasikan array C++, mengisi dan membacanya dengan \texttt{for}, serta menghindari akses di luar batas.
\end{BlokGelap}
\end{column}
\end{columns}
\note{Tanyakan apa yang sudah dibahas di Algoritma minggu ini. Gunakan jawabannya sebagai jembatan ke praktikum.}
\end{frame}

\begin{frame}{Saat pulang nanti, kalian bisa}
\framesubtitle{Target hari ini}
\begin{columns}[T]
\begin{column}{0.47\textwidth}
\begin{Langkah}{1}{Menerapkan}
Menulis program yang menyimpan data sejenis di array satu dimensi dan memprosesnya dengan perulangan: mengisi, menampilkan, menjumlahkan, mencari maksimum, dan menggeser

{\footnotesize\color{ITERAInkMuted}Sub-CPMK 9.1, CPMK0607}
\end{Langkah}
\end{column}
\begin{column}{0.47\textwidth}
\begin{Langkah}{2}{Menjelaskan}
Menelusuri isi array sesudah setiap langkah pemrosesan dan menjelaskan kesalahan indeks, termasuk akses di luar batas

{\footnotesize\color{ITERAInkMuted}Sub-CPMK 9.2, CPMK0608}
\end{Langkah}
\end{column}
\end{columns}
\vspace{8pt}
\begin{Catatan}
Dua kemampuan ini dinilai sama berat di Exit Ticket, tugas, UTS, dan UAS.
\end{Catatan}
\note{Bacakan target dengan pelan. Kata kuncinya menerapkan dan menjelaskan.}
\end{frame}

\begin{frame}{Ingat minggu lalu}
\framesubtitle{Review, 5 menit}
\begin{columns}[T]
\begin{column}{0.31\textwidth}
\begin{BlokGelap}{Pertanyaan 1}
Apa hasil \texttt{int n = 4096} setelah dijumlah digitnya?
\end{BlokGelap}
\end{column}
\begin{column}{0.31\textwidth}
\begin{BlokGelap}{Pertanyaan 2}
Berapa kali badan \texttt{for i 1..3, for j 1..i} dijalankan?
\end{BlokGelap}
\end{column}
\begin{column}{0.31\textwidth}
\begin{BlokGelap}{Pertanyaan 3}
Apa beda soal A dan soal B di UTS?
\end{BlokGelap}
\end{column}
\end{columns}
\vspace{8pt}
Jawab di kertas dulu, lalu cocokkan dengan teman sebelah.\note{Pertemuan pertama sesudah UTS. Bahas singkat kesalahan umum UTS sebelum masuk materi.}
\end{frame}

\Bagian{1}{Masalah pembuka}{Siapa yang di atas rata-rata?}
\note{Masalah pembuka 10 menit.}

\begin{frame}[fragile]{Siapa yang di atas rata-rata?}
\framesubtitle{Masalah minggu ini}
\begin{columns}[T]
\begin{column}{0.55\textwidth}
{\small Dosen ingin tahu berapa mahasiswa yang nilainya di atas rata-rata kelas.

\vspace{6pt}
Dengan pola sentinel minggu 5, rata-rata baru diketahui setelah semua nilai dibaca. Tetapi saat itu nilai-nilai sebelumnya sudah hilang, karena setiap \texttt{cin} menimpa variabel yang sama.\par}
\vspace{6pt}
\begin{Catatan}
\textbf{Diskusi 2 menit.} Apa yang perlu dilakukan supaya semua nilai masih ada saat rata-rata sudah diketahui? Berapa variabel yang dibutuhkan untuk 40 mahasiswa?
\end{Catatan}
\end{column}
\begin{column}{0.41\textwidth}
\begin{Keluaran}[]
Nilai: 70 85 60 90 75 80
Rata-rata   : 76.67
Di atas rata: 3
\end{Keluaran}
\end{column}
\end{columns}
\note{Tampung beberapa jawaban tanpa menilai benar salah.}
\end{frame}

\begin{frame}{Dari langkah ke instruksi}
\framesubtitle{Sambungan dengan Algoritma}
\begin{columns}[T]
\begin{column}{0.31\textwidth}
\begin{Langkah}{1}{Simpan semua}
Semua nilai perlu disimpan, bukan hanya totalnya.
\end{Langkah}
\end{column}
\begin{column}{0.31\textwidth}
\begin{Langkah}{2}{Proses dua kali}
Putaran pertama menghitung rata-rata, putaran kedua membandingkan.
\end{Langkah}
\end{column}
\begin{column}{0.31\textwidth}
\begin{Langkah}{3}{Terjemahkan}
Satu array, dua \texttt{for} dengan indeks 0 sampai n − 1.
\end{Langkah}
\end{column}
\end{columns}
\vspace{10pt}
Langkah 1 dan 2 dilatih di Algoritma. Langkah 3 kita kerjakan hari ini dengan C++.\note{Hubungkan eksplisit ke Algoritma minggu ini.}
\end{frame}

\Bagian{2}{Coba dulu, baru dijelaskan}{25 menit eksperimen di GDB Online}
\note{Struggle 25 menit. Berkeliling, jangan menjelaskan materi dulu.}

\begin{frame}{Misi kalian}
\framesubtitle{Struggle, 25 menit}
\begin{columns}[T]
\begin{column}{0.31\textwidth}
\begin{Langkah}{1}{Deklarasi}
Buat \texttt{int a[5] = \{70, 85, 60, 90, 75\};} dan tampilkan \texttt{a[0]} dan \texttt{a[4]}.
\end{Langkah}
\end{column}
\begin{column}{0.31\textwidth}
\begin{Langkah}{2}{Ulangi}
Tampilkan semua elemen dengan \texttt{for}.
\end{Langkah}
\end{column}
\begin{column}{0.31\textwidth}
\begin{Langkah}{3}{Uji batas}
Tampilkan \texttt{a[5]}. Apa yang muncul?
\end{Langkah}
\end{column}
\end{columns}
\vspace{8pt}
\begin{columns}[T]
\begin{column}{0.47\textwidth}
\begin{Blok}{Boleh}
Bertanya ke teman, mengubah apa saja, sengaja membuat error.
\end{Blok}
\end{column}
\begin{column}{0.47\textwidth}
\begin{BlokAlert}{Belum dulu}
Mencari jawaban jadi di internet atau AI.
\end{BlokAlert}
\end{column}
\end{columns}
\note{Catat kesalahan yang sering muncul untuk dibahas di debrief.}
\end{frame}

\begin{frame}{Apa yang kalian temukan?}
\framesubtitle{Debrief struggle}
\begin{columns}[T]
\begin{column}{0.31\textwidth}
\begin{BlokGelap}{Pertanyaan 1}
Mengapa indeks pertama 0, bukan 1?
\end{BlokGelap}
\end{column}
\begin{column}{0.31\textwidth}
\begin{BlokGelap}{Pertanyaan 2}
Apa yang tampil untuk \texttt{a[5]}, dan apakah sama di setiap komputer?
\end{BlokGelap}
\end{column}
\begin{column}{0.31\textwidth}
\begin{BlokGelap}{Pertanyaan 3}
Bagaimana kalian menghitung jumlah semua elemen?
\end{BlokGelap}
\end{column}
\end{columns}
\vspace{10pt}
Jawaban kalian akan kita uji di bagian berikutnya.\note{Tulis jawaban di papan; buktikan atau patahkan di bagian materi.}
\end{frame}

\Bagian{3}{Banyak nilai, satu nama}{Deklarasi, indeks, perulangan, dan algoritma dasar array}
\note{Materi 40 menit. Tunjukkan setiap contoh langsung di GDB Online.}

\begin{frame}[fragile]{Deklarasi dan indeks}
\framesubtitle{Array satu dimensi}
\begin{columns}[T]
\begin{column}{0.55\textwidth}
\begin{Kode}[]{array.cpp}
int nilai[5] = {70, 85, 60, 90, 75};
cout << nilai[0] << " " << nilai[4] << endl;
nilai[2] = 65;
nilai[3] += 5;
for (int i = 0; i < 5; i++) {
    cout << nilai[i] << " ";
}
cout << endl;
\end{Kode}
\end{column}
\begin{column}{0.41\textwidth}
\only<1>{\begin{TebakDulu}
{\ITERAKickerFont\fontsize{10pt}{12pt}\selectfont\color{ITERAAlert}TEBAK DULU}\par\vspace{4pt}
Sebelum dijalankan, tulis keluarannya di kertas.
\end{TebakDulu}
}\only<2>{\KeluaranFile[]{keluaran/P09/k001.txt}
\begin{Catatan}
Array berukuran 5 punya indeks 0 sampai 4. Ukuran harus konstanta.
\end{Catatan}
}
\end{column}
\end{columns}
\note<1>{Minta mahasiswa menebak keluaran sebelum membuka slide berikutnya. Tanyakan satu atau dua tebakan yang berbeda, lalu buktikan.}
\end{frame}

\begin{frame}[fragile]{Mengisi array dari masukan}
\framesubtitle{Pola baca-proses}
\begin{columns}[T]
\begin{column}{0.55\textwidth}
\begin{Kode}[]{baca.cpp}
int n, a[100];
cin >> n;
for (int i = 0; i < n; i++) {
    cin >> a[i];
}
int total = 0;
for (int i = 0; i < n; i++) {
    total += a[i];
}
cout << "Total " << total << endl;
\end{Kode}
\end{column}
\begin{column}{0.41\textwidth}
\begin{Keluaran}[title={Masukan}]
4
10 20 30 40
\end{Keluaran}
\only<1>{\begin{TebakDulu}
{\ITERAKickerFont\fontsize{10pt}{12pt}\selectfont\color{ITERAAlert}TEBAK DULU}\par\vspace{4pt}
Sebelum dijalankan, tulis keluarannya di kertas.
\end{TebakDulu}
}\only<2>{\KeluaranFile[]{keluaran/P09/k002.txt}
\begin{Catatan}
Deklarasikan array cukup besar (misalnya 100), lalu pakai hanya n elemen pertama.
\end{Catatan}
}
\end{column}
\end{columns}
\note<1>{Minta mahasiswa menebak keluaran sebelum membuka slide berikutnya. Tanyakan satu atau dua tebakan yang berbeda, lalu buktikan.}
\end{frame}

\begin{frame}[fragile]{Mencari maksimum dan posisinya}
\framesubtitle{Algoritma dasar}
\begin{columns}[T]
\begin{column}{0.60\textwidth}
\begin{Kode}[listing options={style=ITERACodeKecil}]{maks.cpp}
int a[6] = {70, 85, 60, 90, 75, 80};
int iMaks = 0;
for (int i = 1; i < 6; i++) {
    if (a[i] > a[iMaks]) {
        iMaks = i;
    }
}
cout << a[iMaks] << " di indeks " << iMaks << endl;
\end{Kode}
\end{column}
\begin{column}{0.36\textwidth}
\only<1>{\begin{TebakDulu}
{\ITERAKickerFont\fontsize{10pt}{12pt}\selectfont\color{ITERAAlert}TEBAK DULU}\par\vspace{4pt}
Sebelum dijalankan, tulis keluarannya di kertas.
\end{TebakDulu}
}\only<2>{\KeluaranFile[listing options={style=ITERAPlainKecil}]{keluaran/P09/k003.txt}
\begin{Catatan}
Simpan indeks, bukan nilainya. Dari indeks, nilai selalu bisa diambil kembali.
\end{Catatan}
}
\end{column}
\end{columns}
\note<1>{Minta mahasiswa menebak keluaran sebelum membuka slide berikutnya. Tanyakan satu atau dua tebakan yang berbeda, lalu buktikan.}
\end{frame}

\begin{frame}[fragile]{Dua putaran: rata-rata lalu bandingkan}
\framesubtitle{Mengapa perlu array}
\begin{columns}[T]
\begin{column}{0.55\textwidth}
\begin{Kode}[]{diatas.cpp}
int a[6] = {70, 85, 60, 90, 75, 80};
int total = 0;
for (int i = 0; i < 6; i++) total += a[i];
double rata = total / 6.0;
int diAtas = 0;
for (int i = 0; i < 6; i++) {
    if (a[i] > rata) diAtas++;
}
cout << rata << " " << diAtas << endl;
\end{Kode}
\end{column}
\begin{column}{0.41\textwidth}
\only<1>{\begin{TebakDulu}
{\ITERAKickerFont\fontsize{10pt}{12pt}\selectfont\color{ITERAAlert}TEBAK DULU}\par\vspace{4pt}
Sebelum dijalankan, tulis keluarannya di kertas.
\end{TebakDulu}
}\only<2>{\KeluaranFile[]{keluaran/P09/k004.txt}
}
\end{column}
\end{columns}
\note<1>{Minta mahasiswa menebak keluaran sebelum membuka slide berikutnya. Tanyakan satu atau dua tebakan yang berbeda, lalu buktikan.}
\end{frame}

\begin{frame}[fragile]{Di luar batas}
\framesubtitle{Kesalahan paling berbahaya}
\begin{columns}[T]
\begin{column}{0.55\textwidth}
\begin{Kode}[]{batas.cpp}
int a[3] = {1, 2, 3};
int jumlah = 0;
for (int i = 0; i < 3; i++) {
    jumlah += a[i];
}
cout << jumlah << endl;
// for (int i = 0; i <= 3; i++) membaca a[3]:
// di luar array, hasilnya tidak bisa ditebak
\end{Kode}
\end{column}
\begin{column}{0.41\textwidth}
\only<1>{\begin{TebakDulu}
{\ITERAKickerFont\fontsize{10pt}{12pt}\selectfont\color{ITERAAlert}TEBAK DULU}\par\vspace{4pt}
Sebelum dijalankan, tulis keluarannya di kertas.
\end{TebakDulu}
}\only<2>{\KeluaranFile[]{keluaran/P09/k005.txt}
\begin{Catatan}
C++ tidak memeriksa batas indeks. Program bisa tetap jalan dengan hasil acak, atau berhenti mendadak.
\end{Catatan}
}
\end{column}
\end{columns}
\note<1>{Minta mahasiswa menebak keluaran sebelum membuka slide berikutnya. Tanyakan satu atau dua tebakan yang berbeda, lalu buktikan.}
\end{frame}

\begin{frame}{Menggeser dan membalik}
\framesubtitle{Menelusuri isi array}
Rotasi kanan \{10, 20, 30, 40, 50\}: simpan elemen terakhir, geser dari belakang, isi depan.
\vspace{8pt}

\small
\begin{tabular}{@{}>{\raggedright\arraybackslash}p{172pt}>{\raggedright\arraybackslash}p{364pt}>{\raggedright\arraybackslash}p{326pt}@{}}
\kepala{Langkah} & \kepala{Isi array} & \kepala{Keterangan}\\
awal & 10 20 30 40 50 & simpan = 50\\
\barisgenap i = 4 & 10 20 30 40 40 & a[4] = a[3]\\
i = 3 & 10 20 30 30 40 & a[3] = a[2]\\
\barisgenap i = 1 & 10 10 20 30 40 & setelah i = 2 dan i = 1\\
akhir & 50 10 20 30 40 & a[0] = simpan\\
\garisdasar
\end{tabular}
\end{frame}

\begin{frame}[fragile]{Tebak keluarannya}
\framesubtitle{Latihan menjelaskan, 3 menit}
\begin{columns}[T]
\begin{column}{0.56\textwidth}
\begin{Kode}[]{tebak.cpp}
int a[5] = {3, 1, 4, 1, 5};
for (int i = 1; i < 5; i++) {
    a[i] = a[i] + a[i - 1];
}
for (int i = 0; i < 5; i++) {
    cout << a[i] << " ";
}
cout << endl;
\end{Kode}
\end{column}
\begin{column}{0.40\textwidth}
\begin{Catatan}
Tulis keluarannya di kertas, karakter demi karakter. Jangan dijalankan dulu.
\end{Catatan}
\end{column}
\end{columns}
\note{Contoh soal Bagian B. Beri waktu 3 menit sebelum slide jawaban.}
\end{frame}

\begin{frame}[fragile]{Tebak keluarannya: jawaban}
\framesubtitle{Latihan menjelaskan}
\begin{columns}[T]
\begin{column}{0.56\textwidth}
\begin{Kode}[]{tebak.cpp}
int a[5] = {3, 1, 4, 1, 5};
for (int i = 1; i < 5; i++) {
    a[i] = a[i] + a[i - 1];
}
for (int i = 0; i < 5; i++) {
    cout << a[i] << " ";
}
cout << endl;
\end{Kode}
\end{column}
\begin{column}{0.40\textwidth}
\begin{Keluaran}[]
3 4 8 9 14 
\end{Keluaran}
\begin{Catatan}
Setiap elemen ditambah elemen sebelumnya yang sudah diperbarui, sehingga terbentuk jumlah kumulatif: 3, 4, 8, 9, 14.
\end{Catatan}
\end{column}
\end{columns}
\note{Tanyakan jawaban yang paling sering salah, lalu telusuri bersama.}
\end{frame}

\begin{frame}{Error yang sering muncul minggu ini}
\framesubtitle{Pesan diambil langsung dari g++}
\footnotesize
\begin{tabular}{@{}>{\raggedright\arraybackslash}p{208pt}>{\raggedright\arraybackslash}p{256pt}>{\raggedright\arraybackslash}p{398pt}@{}}
\kepala{Kesalahan} & \kepala{Contoh} & \kepala{Pesan pertama dari g++}\\
Ukuran negatif & \texttt{int a[-3];} & \texttt{error: narrowing conversion of '-3' from 'int' to 'long unsigned int'}\\
\barisgenap Inisialisasi kelebihan elemen & \texttt{int a[2] = \{1, 2, 3\};} & \texttt{error: too many initializers for 'int [2]'}\\
Indeks bukan bilangan bulat & \texttt{a[1.5]} & \texttt{error: invalid types 'int [3][double]' for array subscript}\\
\barisgenap Menugaskan array ke array & \texttt{b = a;} & \texttt{error: invalid array assignment}\\
Ukuran tidak ditulis & \texttt{int a[];} & \texttt{error: storage size of 'a' isn't known}\\
\garisdasar
\end{tabular}
\vspace{10pt}

{\small Pesan di tabel ini dihasilkan dengan g++ dan opsi -Wall, sama seperti di cek.sh.}
\note{Minta mahasiswa memotret slide ini.}
\end{frame}

\Bagian{4}{Hands-on bertingkat}{45 menit, dari dasar sampai tantangan}
\note{Mahasiswa membuka folder 02 Hands-on/P9 - Array 1D - Hands-on.}

\begin{frame}{Peta hands-on}
\framesubtitle{Folder 02 Hands-on/P9 - Array 1D - Hands-on}
\small
\begin{tabular}{@{}>{\raggedright\arraybackslash}p{36pt}>{\raggedright\arraybackslash}p{267pt}>{\raggedright\arraybackslash}p{110pt}>{\raggedright\arraybackslash}p{83pt}>{\raggedright\arraybackslash}p{341pt}@{}}
\kepala{No} & \kepala{File} & \kepala{Tingkat} & \kepala{Waktu} & \kepala{Yang dilatih}\\
1 & \texttt{handson1-nilai.cpp} & Dasar & 10' & mengisi dan menampilkan, urutan terbalik\\
\barisgenap 2 & \texttt{handson2-statistik.cpp} & Menengah & 15' & rata-rata, maksimum, dua putaran\\
3 & \texttt{handson3-geser.cpp} & Tantangan & 20' & telusuri isi array, perbaiki tiga kesalahan indeks\\
\barisgenap + & \texttt{bonus-frekuensi.cpp} & Pengayaan & sisa & array sebagai pencacah\\
\garisdasar
\end{tabular}
\vspace{10pt}

Kerjakan berurutan. Cocokkan keluaran dengan folder \texttt{expected/}; latihan yang membaca masukan memakai data di folder \texttt{input/}.\note{Saat berkeliling, minta mahasiswa yang mengerjakan latihan tantangan menjelaskan blok PENJELASAN mereka.}
\end{frame}

\begin{frame}{Cara mengecek dan aturannya}
\framesubtitle{Hands-on}
\begin{columns}[T]
\begin{column}{0.47\textwidth}
\begin{Blok}{Di GDB Online}
Salin isi file latihan, lengkapi TODO, klik Run. Bila program membaca masukan, ketik isi file di \texttt{input/}. Bandingkan dengan \texttt{expected/}.
\end{Blok}
\end{column}
\begin{column}{0.47\textwidth}
\begin{BlokGelap}{Di komputer sendiri}
Jalankan \texttt{bash cek.sh}. Skrip mengompilasi dengan \texttt{-Wall -Werror}, memberi masukan dari \texttt{input/}, lalu mencocokkan keluaran.
\end{BlokGelap}
\end{column}
\end{columns}
\vspace{6pt}
\begin{Catatan}
AI boleh untuk memahami pesan error, tidak untuk meminta solusi utuh. Exit Ticket dikerjakan tanpa bantuan apa pun.
\end{Catatan}
\note{Hands-on tidak dinilai langsung; yang dinilai Exit Ticket dan tugas.}
\end{frame}

\Bagian{5}{Exit Ticket dan tugas}{25 menit terakhir, lalu pekerjaan rumah}
\note{Mulai Exit Ticket tepat waktu.}

\begin{frame}{Exit Ticket 9}
\framesubtitle{Individual, 25 menit, tanpa AI dan internet}
\small
\begin{tabular}{@{}>{\raggedright\arraybackslash}p{60pt}>{\raggedright\arraybackslash}p{560pt}>{\raggedright\arraybackslash}p{80pt}>{\raggedright\arraybackslash}p{150pt}@{}}
\kepala{Soal} & \kepala{Isi} & \kepala{Poin} & \kepala{CPMK}\\
A1 & Tulis program yang membaca n nilai ke array dan menampilkan nilai minimum beserta indeksnya & 25 & CPMK0607\\
\barisgenap A2 & Tulis program yang menghitung banyak elemen genap dan ganjil & 25 & CPMK0607\\
B1 & Telusuri isi array setelah sebuah perulangan yang mengubah elemennya & 20 & CPMK0608\\
\barisgenap B2 & Jelaskan mengapa sebuah perulangan menyebabkan akses di luar batas & 20 & CPMK0608\\
B3 & Jelaskan mengapa sebuah masalah butuh array, bukan variabel tunggal & 10 & CPMK0608\\
\garisdasar
\end{tabular}
\vspace{10pt}

{\small Soal A dinilai dari keluaran yang tepat sama. Soal B dinilai dari ketepatan dan alasan. Pelanggaran aturan membuat nilai Exit Ticket ini 0.}
\note{Soal lengkap dibagikan terpisah dan tidak disimpan di repo publik.}
\end{frame}

\begin{frame}{Tugas 9: Sistem Nilai Mahasiswa}
\framesubtitle{Dikumpulkan sebelum Pertemuan 10}
\begin{columns}[T]
\begin{column}{0.47\textwidth}
\begin{Blok}{Bagian A, program, 50 poin}
Buat program dengan ketentuan berikut. Gunakan \texttt{const} untuk ukuran array. Ketentuannya di slide berikut.
\end{Blok}
\end{column}
\begin{column}{0.47\textwidth}
\begin{BlokGelap}{Bagian B, penjelasan, 50 poin}
Komentar maksimal 150 kata dan gambar isi kedua array setelah membaca tiga data contoh. Jelaskan cara menjaga indeks tetap di dalam batas, dan mengapa dua array harus diproses dengan indeks yang sama.
\end{BlokGelap}
\end{column}
\end{columns}
\vspace{8pt}
{\small Data di tugas adalah data kalian sendiri. Rubrik lengkap ada di Modul Belajar 9.}
\note{Ingatkan bahwa penjelasan disalin dari AI mudah dikenali karena tidak menyebut pengalaman mereka sendiri.}
\end{frame}

\begin{frame}{Ketentuan Tugas 9}
\framesubtitle{Sistem Nilai Mahasiswa}
\small
\begin{tabular}{@{}>{\raggedright\arraybackslash}p{87pt}>{\raggedright\arraybackslash}p{788pt}@{}}
\kepala{No} & \kepala{Ketentuan Bagian A}\\
1 & Baca jumlah mahasiswa (maksimal 50)\\
\barisgenap 2 & Baca NIM dan nilai (0 sampai 100) setiap mahasiswa ke dua array sejajar, dengan validasi\\
3 & Tampilkan semua data, nilai tertinggi dan terendah, serta rata-rata kelas\\
\barisgenap 4 & Tampilkan jumlah LULUS (nilai 60 ke atas), TIDAK LULUS, dan persentase kelulusan\\
5 & Fitur pencarian: baca sebuah NIM, tampilkan nilainya atau pesan tidak ditemukan\\
\barisgenap Bonus & Urutkan nilai dari tertinggi ke terendah (materi sorting di Pertemuan 12).\\
\garisdasar
\end{tabular}
\note{Bacakan ketentuan satu per satu dan tanyakan apakah ada yang belum jelas.}
\end{frame}

\begin{frame}{Rubrik Tugas 9}
\framesubtitle{Empat level, sama di semua instrumen}
\footnotesize
\begin{tabular}{@{}>{\raggedright\arraybackslash}p{166pt}>{\raggedright\arraybackslash}p{44pt}>{\raggedright\arraybackslash}p{166pt}>{\raggedright\arraybackslash}p{156pt}>{\raggedright\arraybackslash}p{146pt}>{\raggedright\arraybackslash}p{146pt}@{}}
\kepala{Kriteria} & \kepala{Poin} & \kepala{Sangat baik 85--100} & \kepala{Baik 70--84} & \kepala{Cukup 55--69} & \kepala{Kurang <55}\\
A1 Program berjalan & 20 & tanpa error dan peringatan & ada peringatan & perlu satu perbaikan & tidak terkompilasi\\
\barisgenap A2 Kelengkapan fitur & 15 & semua statistik dan pencarian & satu fitur kurang & dua kurang & tiga atau lebih kurang\\
A3 Validasi dan ketepatan & 15 & validasi tepat, hasil benar untuk n = 1 dan n = 10 & satu salah & dua salah & sebagian besar salah\\
\barisgenap B1 Isi array dan batas indeks & 30 & gambar tepat, alasan jelas & satu langkah keliru & gambar tanpa alasan & tidak ada atau disalin\\
B2 Cerita error dan pengujian & 20 & pesan, sebab, perbaikan, uji & pesan tidak dikutip & hanya perbaikan & tidak ada\\
\garisdasar
\end{tabular}
\vspace{8pt}

{\small Nilai kriteria = poin × skor level. Bagian A total 50, Bagian B total 50.}
\note{Rubrik ini sama strukturnya setiap minggu; yang berubah hanya isi kriteria A2, A3, dan B1.}
\end{frame}

\begin{frame}{Sebelum Pertemuan 10}
\framesubtitle{Persiapan}
\begin{columns}[T]
\begin{column}{0.31\textwidth}
\begin{Langkah}{1}{Selesaikan}
Hands-on yang belum lulus \texttt{cek.sh}, terutama latihan tantangan.
\end{Langkah}
\end{column}
\begin{column}{0.31\textwidth}
\begin{Langkah}{2}{Kumpulkan}
Tugas 9 sebelum Pertemuan 10 dimulai.
\end{Langkah}
\end{column}
\begin{column}{0.31\textwidth}
\begin{Langkah}{3}{Baca}
Modul Belajar 9 untuk mengulang, lalu bagian awal modul berikutnya.
\end{Langkah}
\end{column}
\end{columns}
\vspace{10pt}
Pertemuan 10 membahas array dua dimensi dan string, sejalan dengan Algoritma Pertemuan 10.\note{Tanyakan satu hal yang paling membingungkan hari ini untuk dibuka di review minggu depan.}
\end{frame}

\SlidePenutup{Terima kasih}{Sampai jumpa di Pertemuan 10: array dua dimensi dan string}
\note{Slide penutup.}
\end{document}
